% !TeX program = lualatex
\documentclass[10pt,a4paper]{ltjsarticle}

\usepackage[top=20mm,bottom=22mm,left=20mm,right=20mm]{geometry}
\usepackage{fontspec}
\usepackage{luatexja-fontspec}
\IfFontExistsTF{Hiragino Mincho ProN}{
  \setmainjfont{Hiragino Mincho ProN}
  \setsansjfont{Hiragino Sans}
}{
  \setmainjfont{HaranoAjiMincho}
  \setsansjfont{HaranoAjiGothic}
}
\setmainfont{TeX Gyre Termes}
\setsansfont{TeX Gyre Heros}
\setmonofont{Menlo}[Scale=0.82]

\usepackage{amsmath,amssymb}
\usepackage{booktabs}
\usepackage{tabularx}
\usepackage{array}
\usepackage{multirow}
\usepackage{threeparttable}
\usepackage{float}
\usepackage[section]{placeins}
\usepackage{algorithm}
\usepackage{algpseudocode}
\usepackage{siunitx}
\usepackage{xcolor}
\usepackage{enumitem}
\usepackage{tikz}
\usetikzlibrary{arrows.meta,positioning,fit,calc,shapes.geometric,backgrounds}
\usepackage{fontawesome5}
\usepackage{hyperref}
\usepackage[nameinlink,noabbrev]{cleveref}
\makeatletter
\def\theHALG@line{\thealgorithm.\arabic{ALG@line}}
\makeatother
\floatname{algorithm}{アルゴリズム}
\crefname{algorithm}{アルゴリズム}{アルゴリズム}
\algrenewcommand\algorithmicrequire{\textbf{入力：}}
\algrenewcommand\algorithmicensure{\textbf{出力：}}
\algrenewcommand\algorithmicreturn{\textbf{返す}}
\usepackage[backend=biber,style=numeric,sorting=none,maxbibnames=99]{biblatex}
\addbibresource{references.bib}
\AtBeginBibliography{\footnotesize}

% 印刷と複写を想定し、図表とリンクは無彩色で統一する。
% 図中の区別には、色相ではなく濃淡、線種、形状を用いる。
\definecolor{TorPurple}{gray}{0}
\definecolor{TorLight}{gray}{0.94}
\definecolor{EvidenceBlue}{gray}{0.88}
\definecolor{ResultGreen}{gray}{0.82}
\definecolor{LimitRed}{gray}{0.91}
\definecolor{GrayBox}{gray}{0.96}

\hypersetup{
  colorlinks=true,
  linkcolor=black,
  citecolor=black,
  urlcolor=black,
  pdftitle={Tor CGO暗号状態のメモリフォレンジック},
  pdfauthor={Hiro Nakanishi}
}

\setlength{\parindent}{1em}
\setlength{\parskip}{0.25em}
\setlength{\emergencystretch}{2em}
\setlist{nosep,leftmargin=2em}
\renewcommand{\arraystretch}{1.18}
\newcolumntype{Y}{>{\raggedright\arraybackslash}X}
\newcolumntype{L}[1]{>{\raggedright\arraybackslash}p{#1}}
\newcommand{\code}[1]{\nolinkurl{#1}}
\newcommand{\cgo}{\textsc{CGO}}
\newcommand{\torone}{\textsc{TOR1}}
\newcommand{\pcap}{\textsc{PCAP}}
\newcommand{\centralclaim}[1]{%
  \par\noindent\textbf{本研究の中心主張：}#1\par
}

\title{\vspace{-8mm}
  \textbf{Tor CGO暗号状態のメモリフォレンジック}\par
  \large 回路所有関係の復元とTLS・CGO状態を用いた保存通信の復号}
\author{Hiro Nakanishi\\独立研究者\\\href{mailto:quantumsity@protonmail.com}{quantumsity@protonmail.com}}
\date{}

\begin{document}
\maketitle
\vspace{-5mm}

\begin{abstract}
TorのCounter Galois Onion（CGO）は、ホップごと・方向ごとに暗号状態を更新する。
インシデント調査でTorクライアントのメモリを取得できても、暗号状態らしいバイト列だけでは、取得時点の所有回路への所属、入口からのホップ順、取得時点での利用を立証できない。
本研究は、C Torのコアダンプと対応するDWARF情報を用い、Tor自身の所有リストから回路を確定し、\code{cpath}から入口起点の順序を復元する手法を提案する。
さらに、回路が保持する入口TLS接続をたどり、TLS~1.3秘密、レコード番号、Tor回路ID、CGO状態を結合することで、取得直後に記録された保存PCAP内のリレーセルをTLS層からCGO層までオフラインで復号する。
Linux arm64上の複数のTor 0.4.9系ビルドでは、独立12プロセスの制御下144/144回路と432/432ホップを復元した。
コアダンプ内の状態と取得後のパケットが対応する条件では、Torと隔離網を毎回作り直した10/10独立起動で、所有関係、TLS AEAD、一意なCGO認証を検証し、評価用の正解と一致する493バイトを復元した。通信を一時停止せず、セル・TLS境界を観測しない固定遅延で取得した6/6独立起動でも、異なるTLSレコード番号と宛先ホップについて同じ結果を得た。Control Protocol由来の正解回路IDは候補選択に用いなかった。
Tor ProjectがDebian 12向けに配布するストリップ済みTor 0.4.9.17と、対応する外部デバッグシンボルを用いた3/3独立起動でも、36/36回路と108/108ホップを復元した。同じ配布版を用いた別の6起動を、公式0.4.9.17ソースから構築した再生器で処理し、一意なCGO認証済みリレーセル本体を6/6で復元した。内訳は評価用マーカーを含むDATAセル5件とSENDMEセル1件だった。
一方、現行状態による直前セルの復号は3/3で失敗し、活動中に端点を順次取得した場合は2.30--3.17秒の差で状態一致が0/36となった。
したがって本手法は、端点メモリを取得できる調査者に、回路、経路、暗号方式、取得後の通信内容の対応関係を示すが、過去通信の包括的復号や遠隔からのTor匿名性破壊を与えない。

\noindent\textbf{キーワード：} Tor、Counter Galois Onion、メモリフォレンジック、コアダンプ、TLS、オフライン復号
\end{abstract}

\begin{center}
\textbf{Memory Forensics of Tor CGO Cryptographic State:\\
Recovering Circuit Ownership and Decrypting Stored Traffic with TLS and CGO State}
\end{center}
\begin{quote}\small
\textbf{Abstract.}
Counter Galois Onion (CGO) updates cryptographic state for every circuit hop and direction.
Finding key-like bytes in a Tor client memory image does not establish which live circuit owned them, which hop was first, or whether the state was active at acquisition time.
We present a forensic method that starts from the ownership relationships maintained by C Tor.
The method resolves Tor's circuit-owner list using matching DWARF information, recovers entry-oriented hop order from \texttt{cpath}, and links the entry TLS object to captured TLS records, Tor circuit identifiers, and per-hop CGO state.
Across twelve independent processes covering four Tor 0.4.9 build conditions on Linux arm64, the analyzer recovered 144/144 controlled circuits and 432/432 established CGO hops.
Across three further independent processes using the stripped Tor 0.4.9.17 package and matching external debug symbols distributed by the Tor Project for Debian 12, the analyzer recovered 36/36 circuits and 108/108 established CGO hops.  In six separate packaged-Tor runs, a unique relay body passed CGO authentication in all six; five contained the evaluation marker, while one was a different authenticated relay body.
In ten fresh, independent Tor runs, a circuit list containing only circuits with verified ownership, path structure, and initialized CGO contexts constrained both TLS and CGO extraction.  A client core and packet capture spanning its acquisition point were then sufficient to authenticate and decrypt TLS~1.3 records, enumerate all cells without a Control Protocol truth identifier, select the unique CGO-authenticated relay cell, recognize its destination hop, and recover all 493 non-tag bytes of the relay body.  Six further independent runs succeeded without pausing the application stream or observing a relay-cell or TLS-record boundary before acquisition, across six distinct TLS record sequence numbers and destination hops one and two.
The current state failed to decrypt the immediately preceding cell in 3/3 runs, and sequential endpoint acquisition during traffic preserved circuit links but lost all 36/36 instantaneous state matches.
The contribution is therefore a method for verifying the correspondence among endpoint memory, Tor circuits, and post-acquisition traffic, not a remote cryptanalytic break or general deanonymization attack.

\noindent\textbf{Keywords:} Tor, Counter Galois Onion, memory forensics, core dump, TLS, offline decryption
\end{quote}

\section{はじめに}

Torは、複数のリレーを通る回路と多層の暗号処理によって、通信元と宛先の結び付きを困難にする~\cite{dingledine2004tor}。
この性質は通常の通信保護には有効である。一方、侵害された端末を調べる場面では、調査者が取得したメモリのどこに、どの回路の状態があり、その状態から何を復元できるかが問題になる。

CGOは認証付きのonion encryptionをTorへ導入する方式であり、セル処理に伴って鍵、ノンス、累積状態$T'$を更新する~\cite{torproposal359,degabriele2025cgo}。
更新性は、単純な鍵検索を不十分にする。メモリ上にCGO構造体を発見しても、それが開いた回路に属するのか、どのホップを表すのか、どの時点のセルへ適用できるのかは分からないからである。

本研究は、次の問いに答える。

\begin{quote}
\textbf{研究質問：} 未改変C Torの制御下の継続通信中に取得したコアダンプと関連証拠を用いて、どの条件の下で、CGO/\torone{}の回路状態と通信内容をどの範囲までオフラインで復元できるか。
\end{quote}

\centralclaim{コアダンプに残るTLS書込みレコード番号と、保存PCAP内で認証できる取得後最初のレコード番号が一致する場合、回路所属、入口起点のホップ順、入口TLS、ホップ別CGO状態を対応付け、そのレコード内でCGO認証に成功するリレーセル本文を復元できる。これはCGOそのものを破る攻撃ではなく、端点メモリ取得後に成立するフォレンジック手法である。}

第\ref{sec:contributions}節に貢献、第\ref{sec:background}節にCGOと調査対象、第\ref{sec:related}節に先行研究との差を示す。第\ref{sec:method}節は証拠に基づく提案手法、第\ref{sec:implementation}節は評価実装、第\ref{sec:experiment}節は実験設計、第\ref{sec:results}節は結果を扱う。第\ref{sec:security}節で抽出状態が与える攻撃能力の境界を検証し、第\ref{sec:discussion}節と第\ref{sec:limitations}節に実務上の意味と限界を示す。

解析器、実験スクリプト、機微情報を除いた評価結果は、公開リポジトリで提供する\footnote{実装・補足資料：\url{https://github.com/quantumshiro/tor-cgo-analysis}}。

\section{本研究の貢献}\label{sec:contributions}

本研究の貢献は、次の四点に要約できる。

\begin{enumerate}[label=\textbf{C\arabic*.},leftmargin=3em]
  \item \textbf{回路所有関係を起点とする静的メモリ解析手法。}
  鍵や構造体のパターン走査ではなく、Tor自身の所有リストからクライアント起点回路（origin circuit）を列挙する。\code{origin_circuit_t.cpath}を入口とし、マジック値、リスト添字、双方向リンク、閉路を検査することで、候補の回路所属と入口起点のホップ順を同時に検証する。複数バージョン・最適化条件の独立12プロセスで、制御下144/144回路と432/432ホップを復元した。配布パッケージ版の独立3プロセスでも、36/36回路と108/108ホップを復元した。

  \item \textbf{回路とTLS・CGO状態の対応関係。}
  回路ID、ストリーム宛先、入口チャネル、ホップ別CGO/\torone{}、方向別CGO状態を、Tor構造体間のポインタ関係に沿って対応付ける。これにより、メモリ内の鍵らしい値を発見するだけの手法と異なり、「どの回路の、どのホップの、どの方向の状態か」を説明できる。

  \item \textbf{TLSレコードとCGO層を順に処理する保存通信の復号。}
  所有回路から入口OpenSSLオブジェクトを特定し、TLS~1.3通信秘密と書込みレコード番号を取得する。通信復元の対象は、所有関係、経路構造、暗号コンテキストを検証できた全ホップCGO回路である。PCAPのTLS認証復号、全Torセル候補の列挙、CGO層処理を順に行う。自己ビルド版では、通信を一時停止した10/10独立起動と、ストリームを止めずセル・TLS境界を観測しない固定遅延で取得した6/6起動で、一意なCGO認証済みリレーセル本体を復元した。公式パッケージ版でも6/6起動でCGO認証に成功し、評価用マーカーを含むリレーDATAセルは5/6で得られた。Control Protocolや動的計測による正解情報は復号入力に含めず、候補確定後の評価だけに用いた。

  \item \textbf{成功条件と誤認境界の定量化。}
  構築中の未初期化列挙値、欠損ページ、誤った実行ファイル、誤鍵、取得時点のずれ、直前セル、既知の後方向再同期攻撃を負の対照として評価した。その結果、現行状態はコアダンプ取得後の次セルには使えるが、過去通信の包括的復号、遠隔攻撃、一般的な匿名性破壊を意味しないことを示した。
\end{enumerate}

以上の貢献はCGOの新しい暗号解読ではない。CGO固有の貢献は、ホップごとに更新される認証付き状態列を回路と時点へ帰属させ、保存通信へ正しい順序で適用する点にある。より一般化できる貢献は、所有構造を起点に、型情報、トランスポート層の認証状態、上位暗号状態を段階的に結合し、各検査の成功・失敗・未評価を残す方法である。

\section{背景}\label{sec:background}

\subsection{Tor回路とリレーセル}

Torクライアントは通常、入口、中継、出口の順にリレーを選び、回路を構築する。
クライアントから送る前方向（forward）セルには、出口側から入口側へ暗号層を重ねる。受信する後方向（backward）セルは逆の方向に処理される。
セルは入口とのTLS接続で運ばれるため、ネットワーク上で観測したパケットは、TLS層とonion層の双方で保護されている。

C Torのorigin circuitは、回路ID、入口チャネル、ストリーム、経路を保持する~\cite{ctorsource04916}。
経路は\code{crypt_path_t}の循環双方向リストである。
重要なのは、リスト内の最小アドレスが入口なのではなく、\code{origin_circuit_t.cpath}が入口ホップを指す点である。
したがってアドレス順の推測は不要であり、所有構造の意味論から向きを決められる。

\subsection{CGOの状態と時間方向}

Proposal~359は、リレーセルの509バイトを処理するCGOを定義する~\cite{torproposal359}。
各ホップ・各方向は、疑似乱数関数用の鍵$K$、ノンス$N$、認証に用いる累積値$T'$などを保持する。
CGOはセルを処理するたびに状態を更新する。概念的には、状態$S_i$とセル$C_i$から出力と次状態を得る。

CGOはC Tor 0.4.9.5でクライアントとリレーが交渉できる安定版機能として公開されたが、Proposal~359自体の状態は本稿執筆時点でOpenである~\cite{torstable0495,torproposal359}。また0.4.9.3-alphaから0.4.9.11には、最終ホップ以外で輻輳制御を有効にせず、その結果CGOを交渉しない実装不具合があった。0.4.9.12は、対応する全ホップで輻輳制御とCGOを交渉するよう修正した~\cite{torrelease04912}。したがって、方式の混在はリレー能力だけでなく、クライアント版と回路構築条件にも依存する。列挙値だけでは確立済み暗号方式を判定できないため、取得時点におけるホップ別暗号コンテキストの初期化状態も検証対象とする。

\begin{equation}
  (C'_i, S_{i+1}) = \operatorname{CGO}(C_i,S_i).
\end{equation}

スナップショットが保持するのは取得時点の$S_t$である。$S_t$から、同じ状態を起点とする将来セルを逐次処理できる。一方、更新前の$S_{t-1}$が残っていなければ、直前セル$C_{t-1}$を復号できない。この時間方向は、一般的なディスク復号鍵の回収とは異なる。

\begin{figure}[t]
\centering
\begin{tikzpicture}[
  x=1mm,y=1mm,>=Latex,
  state/.style={draw,rounded corners,minimum width=26mm,minimum height=11mm,align=center,font=\small},
  cell/.style={draw,minimum width=22mm,minimum height=10mm,align=center,font=\small},
  note/.style={font=\footnotesize,align=center}
]
\node[state,fill=GrayBox] (s0) at (0,0) {$S_{t-1}$\\更新前状態};
\node[cell,fill=LimitRed] (past) at (35,0) {直前セル\\$C_{t-1}$};
\node[state,fill=TorLight] (st) at (70,0) {$S_t$\\取得した状態};
\node[cell,fill=EvidenceBlue] (next) at (105,0) {次セル\\$C_t$};
\node[state,fill=ResultGreen] (s1) at (140,0) {$S_{t+1}$\\更新後状態};
\draw[->,thick,shorten <=1mm,shorten >=1mm] (s0) -- (past);
\draw[->,thick,shorten <=1mm,shorten >=1mm] (past) -- node[above=1mm,note]{処理} (st);
\draw[->,thick,shorten <=1mm,shorten >=1mm] (st) -- node[above=1mm,note]{復号可能} (next);
\draw[->,thick,shorten <=1mm,shorten >=1mm] (next) -- node[above=1mm,note]{処理} (s1);
\coordinate (rollback) at ($(st.south)+(0,-15mm)$);
\draw[->,dashed,black!75,shorten >=1mm] (st.south) -- (rollback) -| (past.south);
\coordinate (rollbacklabel) at ($(past.south |- rollback)!0.5!(rollback)$);
\node[note,text=black!75,below=3mm of rollbacklabel] {現行状態から更新前状態は復元できない};
\node[note,above=7mm of st] {コアダンプ取得時点};
\draw[densely dotted] (st.north) -- ++(0,9mm);
\end{tikzpicture}
\caption{CGO状態の時間境界。現行状態は同期した次セルには使えるが、更新前状態を失った直前セルには一般に使えない。}
\label{fig:time}
\end{figure}


Tor Control Protocolは、稼働中のTorへ問い合わせ、回路ID、経路、ストリーム状態などを取得できる~\cite{torcontrol}。
しかし、ホップ別のCGO/\torone{}実体、OpenSSL内部のTLS通信秘密、方向別CGO状態は公開しない。
また、プロセス終了後のコアダンプには接続できない。
Control Protocolの記録は評価用の正解に限定し、解析入力から除外する。

\section{先行研究}\label{sec:related}

\subsection{Onion encryptionとCGO}

初期のTorは、低遅延のオニオンルーティングを実用化した一方、リレー暗号の改変耐性には形式的な課題があった~\cite{dingledine2004tor}。
DegabrieleとStamはonion encryptionを形式化し、Torのrelay cryptographyを安全性定義の観点から分析した~\cite{degabriele2018untagging}。
その後、Degabrieleらは、認証と高速処理を両立するCGOを提案し、構成と性能を示した~\cite{degabriele2025cgo}。
さらに、leaky pipes、関連データ、部分的前方秘匿性を含む安全性モデルが検討された~\cite{degabriele2025secure}。

これらの研究は暗号方式の安全性と状態遷移を扱うが、取得した実プロセスメモリ内で、状態を取得時点の所有回路へ帰属させる方法は扱わない。
Proposal~359も実装が保持すべき状態を定義するが、フォレンジック取得、ページ欠損、ABI差、TLSとの証拠結合は対象外である~\cite{torproposal359}。

\subsection{Tor・暗号鍵・TLSのメモリフォレンジック}

AlfosailとNorrisは、Torクライアントのメモリ内に残る痕跡を対象とする調査手順を提案した~\cite{alfosail2021tor}。
MuirらはTor Browser Bundleを監査し、Windowsホストに残る利用痕跡を評価した~\cite{muir2019forensic}。
ArshadらはWindows~10とAndroid~10で、レジストリ、ストレージ、メモリ、Zramに残る痕跡を比較した~\cite{arshad2021forensic}。
これらはTor利用の立証や端末アーティファクトの収集に有用である。ただし、C Tor内部の回路所有関係、入口起点の暗号経路、ホップ別CGO状態、TLSからCGOまでの連続した復号は研究対象としていない。

Tor Projectは、メモリ中の統計値がプライバシーへ与える影響も分析している~\cite{tor2017memory}。
これはメモリ内情報を最小化する設計上の研究であり、クライアント回路の暗号状態を取得後に解釈する本研究とは目的が異なる。


暗号鍵をRAMから探す研究では、鍵スケジュールの構造を手掛かりにAESなどの鍵を識別できる~\cite{maartmannmoe2009persistence}。
Cold boot研究は、劣化したRAMから鍵を再構成できることを示した~\cite{halderman2008coldboot}。
これらの成果は、暗号状態がメモリ証拠になり得ることを示す。
しかし、鍵が見つかっただけでは、Torのどの回路、どのホップ、どの方向がその鍵を所有するか分からない。
本研究は鍵探索を起点にせず、Torの所有関係を起点に暗号状態へ到達する。


X-Ray-TLSは、プロセスメモリからTLS~1.2、TLS~1.3、QUICのセッション秘密を抽出し、通信を透過的に復号する~\cite{moriconi2024xraytls}。Baierらは、TLS鍵素材のメモリ内識別・抽出手法を体系化し、実装差、TLS~1.3、取得後解析の課題を整理した~\cite{baier2024tlsmemory}。さらに、規格固有文字列から鍵導出関数を見つけ、主要TLS実装へフックする方法も提案されている~\cite{baier2025allyourtlskeys}。これらは、本研究がTLS秘密の抽出自体を新規性として主張できないことを示す重要な隣接研究である。

しかし、入口リンクのTLS平文を得ても、その中のリレーセルはホップ別のCGO層で保護されている。また、TLS秘密だけからは、Torのどの所有回路、入口起点のどのホップ、どの取得時点のCGO状態を適用すべきか決まらない。本研究の差分は、TLS秘密抽出ではなく、Tor構造体間のポインタ関係を起点としてTLSオブジェクト、リンク回路ID、ホップ別CGO状態、取得時点を対応付ける点にある。表~\ref{tab:related-matrix}に調査上の差を示す。

\begin{table}[!tbp]
\centering
\caption{関連手法が解決する調査上の範囲}
\label{tab:related-matrix}
\small
\begin{tabularx}{\linewidth}{Yccccc}
\toprule
手法 & Tor所有回路 & 入口順 & TLS復号 & CGO復号 & TLSレコード番号とCGO状態の対応 \\
\midrule
Tor端末痕跡研究~\cite{alfosail2021tor,muir2019forensic} & -- & -- & -- & -- & -- \\
一般鍵探索~\cite{maartmannmoe2009persistence} & -- & -- & -- & -- & -- \\
X-Ray-TLS~\cite{moriconi2024xraytls} & -- & -- & $\checkmark$ & -- & -- \\
TLS鍵抽出研究群・体系化研究~\cite{baier2024tlsmemory,baier2025allyourtlskeys} & -- & -- & $\checkmark$ & -- & -- \\
本手法 & $\checkmark$ & $\checkmark$ & $\checkmark$ & $\checkmark$ & $\checkmark$ \\
\bottomrule
\end{tabularx}
\end{table}

実務では、本手法より軽い方法で調査目的を満たせる場合がある。DFRWS EU 2026の演習は、端末内のlocalhostでTorへ入る前の通信を取得し、暗号化されたTorトンネルの内側に相当する情報を解析する課題を扱っている~\cite{dfrws2026localtor}。これは事前にloopbackを取得できる場合には有力であり、アプリケーション内容だけが目的なら本手法より単純である。一方、どのTor回路・ホップ・暗号方式がその通信を運んだか、また取得後のコアから何を立証できるかは別の問題である。表~\ref{tab:practical-alternatives}に、現場で競合する取得方法との使い分けを示す。

\begin{table}[!tbp]
\centering
\caption{実務的な代替手法との比較}
\label{tab:practical-alternatives}
\small
\begin{tabularx}{\linewidth}{L{28mm}Y Y L{25mm}}
\toprule
手法 & 得られる主な情報 & 得られない主な情報 & 適する場面 \\
\midrule
通常ログ・Control Protocol & 稼働中の回路ID、経路、ストリーム状態 & ホップ別暗号コンテキスト、TLS秘密、終了後の再解析 & 生存プロセスを即時照会できる \\
localhost／SOCKS側PCAP~\cite{dfrws2026localtor} & Torへ入る前のアプリケーション側通信 & onion回路への帰属、ホップ方式、入口リンク上のセル & 事前取得があり内容確認が目的 \\
汎用TLS鍵抽出~\cite{moriconi2024xraytls,baier2024tlsmemory} & 入口TLS内のTorセル & CGO層のリレーセル本体、所有回路、入口順 & TLS内部だけを確認する \\
コアダンプ単独の本手法 & 取得時点の所有回路、入口順、CGO/\torone{}状態 & 通信本文 & コアダンプだけが残る事後調査 \\
コアダンプとPCAPを用いる本手法 & 上記に加え、取得時点との対応を検証したTLSセルとCGO認証済みリレーセル本体 & 過去の任意通信、上位層暗号の自動解除 & コアダンプと入口PCAPが同じ取得時点に対応する \\
\bottomrule
\end{tabularx}
\end{table}

\subsection{未解決だった調査上の差分}

先行研究を整理すると、未解決の部分は「暗号状態の発見」ではなく「状態の帰属と時点」である。
Control Protocolは生きているプロセスの回路を説明できるが、暗号内部状態や取得後解析を提供しない。
Torフォレンジック研究は端末痕跡を扱うが、CGOの状態意味論を扱わない。
暗号鍵・TLSフォレンジックは鍵やTLS平文を得るが、Torの回路所有関係とCGO層を復元しない。

本研究の新規性は、この三つを検証可能な対応関係として示した点にある。すなわち、Tor内部の所有関係から対象回路を確定し、入口からホップを並べ、同じ回路が保持するTLS状態とCGO状態を保存PCAPへ適用する。この組合せにより、構造らしい残存断片と取得時点の所有状態を区別し、リレー本文まで到達できる。

\section{提案手法}\label{sec:method}

\subsection{手法の全体像と解析問題}

本手法の要点は、メモリ全体から暗号鍵を探すことではない。Torが管理している所有関係を順方向にたどり、最後に暗号状態へ到達することである。解析順序は、\textbf{所有回路の確定}、\textbf{入口とホップ順の確定}、\textbf{暗号方式と初期化状態の確定}、\textbf{TLSおよびCGO秘密の抽出}、\textbf{保存通信との照合}である。

この順序には理由がある。ヒープ上には、終了済み回路の残存構造、構築中のゼロ初期化領域、解放済みの鍵、別回路のOpenSSLコンテキストが残る可能性がある。暗号値だけを先に発見すると、それがどの回路に属するかを後から証明できない。本手法では、Tor自身が維持する\code{global_origin_circuit_list}を探索の起点にする。そこから到達できない構造体や鍵は、値が暗号学的にもっともらしくても、取得時点の所有回路へ帰属できない。

入力にPCAPが含まれるかによって、到達できる解析結果が異なる。コアダンプ、対応する実行ファイル、DWARFだけからは、取得時点の所有回路、ライフサイクル、入口起点の経路、ストリーム宛先、ホップ別CGO/\torone{}、方向別CGO状態を復元する。対象プロセスは終了済みでもよい。

さらに、コアダンプ取得時点をまたいで入口リンクを記録したPCAPがある場合は、コア内のTLS書込みレコード番号に対応するレコードをAEAD認証し、その中のTorセルへCGO状態を適用する。この通信解析は、コアダンプを取得した後にTorを再開できる場合に適する。クラッシュ時のコアダンプしかなく、その後の通信が存在しない場合は、回路状態の解析までとなる。更新前状態の別コピーが残っていない限り、現行CGO状態からコアダンプ取得前のセルへ戻れないためである。


\paragraph{脅威モデルと信頼境界。}
調査者は対象端末を取得する正当な権限を持ち、取得器、解析器、ハッシュ計算環境を信頼できるものとする。コアダンプとPCAPは取得後に改変されず、候補ELF/DWARFはコア写像との互換性検査を通り、OpenSSL内部プロファイルは対象ビルドで検証済みであることを前提とする。対象プロセスが侵害されていても、残存構造、未初期化領域、部分欠損は検査対象に含める。一方、取得前に任意メモリ書換えが可能な攻撃者が、Torの所有リスト、回路構造、暗号状態、PCAPを相互整合するよう意図的に偽造するアンチフォレンジックは対象外である。

所有リストへの掲載と構造検証が示すのは、取得物内部でTorの所有関係と状態が整合することである。これは取得手続そのものの法的な証拠保全を代替しない。TLSおよびCGOの認証タグの検証は、抽出した秘密の下で対象バイト列が認証されたことを示すが、その通信を特定人物が意図したことや、端末外の送信元帰属までは示さない。結果記録には、取得時点、入力ハッシュ、各検査の成功・失敗・未評価、およびその理由を含める。

本手法への必須入力は、時点$t$に取得したTorプロセスのコアダンプ$M_t$と、$M_t$を生成した実行ファイルに対応する型情報$B$である。コアダンプは、プロセスの仮想メモリとレジスタを保存した取得後解析用のスナップショットである。候補ELFのBuild IDを記録し、コアの\code{NT_FILE}対応表が示す主実行ファイル領域と候補ELFの対応バイト列が、第~\ref{sec:implementation}節で定める75\%以上一致した場合だけ$B$を採用する。保存通信を復号する場合だけ、時点$t$をまたいでクライアントと入口リレーの間を記録したパケットキャプチャ$P$を追加する。

解析単位はorigin circuitである。回路$c$について、本手法は次の結果記録$R_c$を生成する。

\begin{equation}
R_c=(I_c,L_c,H_c,A_c,Q_c,U_c,Z_c,V_c).
\label{eq:result-record}
\end{equation}

$I_c$はTor内部の回路識別子と入口リンク上の回路ID、$L_c$は構築中・確立・終了予約済みなどのライフサイクル、$H_c=(h_0,\ldots,h_{n-1})$は入口から並べたホップ列である。$A_c$は回路が所有するストリームと接続先、$Q_c$は各ホップ・各方向の確立済み暗号方式とCGO累積状態の更新有無、$U_c$はTLSレコードとCGO状態がコアダンプ取得時点に対応する情報、$Z_c$は復元できたTorセルまたはリレーセル本体である。$V_c$は、所有リストの完全列挙、回路構造、暗号コンテキスト、TLS AEAD、CGO認証の各検査を、成功・失敗・未評価の三値と理由で記録する。入力不足または検証失敗があるフィールドは未評価とし、その理由を記録する。

Proposal~359で定義されるCGOの一方向状態を、実装上のアドレスではなく次の論理値として表す~\cite{torproposal359}。

\begin{equation}
S_{c,d,h}(t)=(K_J,H_J,K_S,H_S,N,T').
\label{eq:cgo-state}
\end{equation}

$c$は回路、$d$は前方向（forward）または後方向（backward）、$h$はホップ、$t$は取得時点である。$K_J$と$K_S$は二つのAES-128鍵、$H_J$と$H_S$は二つのPOLYVAL鍵、$N$はノンス、$T'$はセル処理で更新される累積状態である。各成分は16バイトであり、一方向につき合計96バイトとなる。解析に必要なのは、この論理値を回路・方向・ホップへ帰属させることである。回路帰属と入口順序の判定はメモリ内のアドレス順や固定オフセットに依存しない。一方、OpenSSL内部から秘密値を抽出する処理は、後述するバージョン別プロファイルに依存する。

\subsection{所有関係から通信復元へ進む処理手順}

手法を、コアダンプだけを用いる回路状態復元（アルゴリズム~\ref{alg:core-recovery}）と、保存PCAPを加える通信復元（アルゴリズム~\ref{alg:traffic-recovery}）に分けて示す。各処理の成立には、直前までの検証がすべて成功している必要がある。後段の検証成功は、前段の未成立を補わない。

所有リスト$O_t$を完全に読み出し、その位置$j$にある要素$c$を解析対象とする。構造体の所属情報を検証する条件を次式で表す。
\begin{equation}
\operatorname{Structured}(c,j)=
[\operatorname{idx}(c)=j]\land[\operatorname{magic}(c)=M_{\mathrm{origin}}].
\label{eq:structured}
\end{equation}

アルゴリズム中の\textsc{VerifyExecutable}は、候補ELFのBuild IDを記録し、コアダンプの\code{NT_FILE}が示す主実行ファイル領域とELFの対応バイト列を照合する。\textsc{SafeRead}は、対象範囲が読み取り可能である場合だけ型付き値を返し、それ以外は失敗理由を返す。\textsc{RecoverOwnedMetadata}は、この検査後に回路識別子、ライフサイクル、所有ストリームを復元する。\textsc{RecoverCryptoState}は各ホップの方式識別子、前方向と後方向の暗号コンテキスト、式~\eqref{eq:cgo-state}のCGO状態を復元する。Proposal~359では$T'$をゼロで初期化し、その後のセル処理で更新するため、$T'\neq0$は、その方向で$T'$が初期値から変化したこととして別に記録する~\cite{torproposal359,ctorsource04916}。

\textsc{AuthenticateExactSequence}は、コアダンプに保存された書込みレコード番号だけを用い、TLSの関連データ付き認証暗号（Authenticated Encryption with Associated Data: AEAD）を検証する。\textsc{EnumerateRelayCells}は、その認証済みレコードに含まれる対象回路の全セルを返す。Control Protocol由来の正解回路IDで一件へ絞らない。\textsc{ApplyCGOFromSnapshot}は各候補を同じ取得時点状態の複製へ入口順に適用し、CGO認証タグを認識したホップと、タグを除く493バイトのリレー本文を返す~\cite{torproposal359}。
所有リストの完全列挙と、個別回路の構造検証は別の結果として記録する。これにより、回路構造が破損していても、そのポインタが取得時点の所有リストに含まれていたという観測事実を残せる。
入口からのホップ列は、$h_0=c.\mathrm{cpath}$、$h_{i+1}=\operatorname{next}(h_i)$として復元する。各ホップで前後リンクの相互一致を検査し、列が$h_0$へ戻る場合だけ閉路を受理する。

保存PCAPを処理する場合、TLS~1.3の通信秘密を$S$、レコード番号を$i$とし、RFC~8446に従って鍵$k$、初期化ベクトル$v$、ノンス$n_i$を導出する~\cite{rfc8446}。ハッシュ$\mathcal H$と鍵長$\ell$は暗号スイートから決まり、本評価の\code{TLS_AES_256_GCM_SHA384}では$\mathcal H=\mathrm{SHA384}$、$\ell=32$である。
\begin{align}
  k &= \operatorname{HKDF\mbox{-}Expand\mbox{-}Label}_{\mathcal H}(S,\mathrm{key},\ell),\label{eq:tls-key}\\
  v &= \operatorname{HKDF\mbox{-}Expand\mbox{-}Label}_{\mathcal H}(S,\mathrm{iv},12),\\
  n_i &= v \oplus \operatorname{I2OSP}_{12}(i).
\end{align}

\begin{algorithm}[H]
\caption{所有関係を起点とする回路状態復元}
\label{alg:core-recovery}
\small
\begin{algorithmic}[1]
\Require コアダンプ$M_t$、候補実行ファイルと型情報$B$
\Ensure 回路ごとの結果記録集合$\mathcal{R}=\{R_c\}$
\State $\mathcal{R}\gets\emptyset$
\If{$\Call{VerifyExecutable}{M_t,B}$が不成立}
  \State \Return $\{\operatorname{error}(\text{実行ファイル不一致})\}$
\EndIf
\State $\Lambda\gets\Call{ResolveTypes}{B}$
\State $O_t\gets\Call{ReadOwnerList}{M_t,\Lambda}$
\If{$O_t$の領域、要素数、容量、全要素を検証できない}
  \State \Return $\{\operatorname{incomplete}(\text{所有リスト取得不能})\}$
\EndIf
\For{$j\gets0$ \textbf{to} $|O_t|-1$}
  \State $(I_c,L_c,H_c,A_c,Q_c,U_c,Z_c,V_c)\gets(\bot,\bot,\bot,\bot,\bot,\bot,\bot,\emptyset)$
  \State $c\gets O_t[j]$; 所有リストへの掲載を記録
  \If{$\Call{SafeRead}{M_t,c,\Lambda}$が失敗}
    \State $V_c$へ回路構造の失敗理由を記録
    \State $R_c\gets(I_c,L_c,H_c,A_c,Q_c,U_c,Z_c,V_c)$; $\mathcal{R}\gets\mathcal{R}\cup\{R_c\}$; \textbf{次の要素へ}
  \EndIf
  \If{$\operatorname{Structured}(c,j)$}
    \State $(I_c,L_c,A_c)\gets\Call{RecoverOwnedMetadata}{c}$
    \State $H_c\gets\Call{WalkFromEntry}{c.\mathrm{cpath}}$
  \EndIf
  \If{$\operatorname{Structured}(c,j)$かつ$H_c$が双方向リンクと閉路の検査を通過}
    \State 回路構造と入口起点経路の検証成功を記録
    \State $Q_c\gets\Call{RecoverCryptoState}{c,H_c}$
    \If{回路が活動中かつ確立状態で、全ホップの方式別暗号コンテキストが整合}
      \State 暗号コンテキストの初期化確認成功を記録
    \EndIf
  \EndIf
  \State 未実行の後段検査を$V_c$で未評価とし、理由を記録
  \State $R_c\gets(I_c,L_c,H_c,A_c,Q_c,U_c,Z_c,V_c)$
  \State $\mathcal{R}\gets\mathcal{R}\cup\{R_c\}$
\EndFor
\State \Return $\mathcal{R}$
\end{algorithmic}
\end{algorithm}

前方向については、入口から出口までの全ホップ状態を
\begin{equation}
\mathcal{S}_{c,\mathrm{forward}}(t)=
\bigl(S_{c,\mathrm{forward},0}(t),\ldots,
S_{c,\mathrm{forward},n-1}(t)\bigr)
\label{eq:cgo-state-sequence}
\end{equation}
と表す。候補セルごとにこの状態列を複製し、Torの前方向処理と同じ入口順で各層へ適用する。あるホップでCGO認証に成功した場合に、そのホップを宛先として返す。したがって、以下の\textsc{ApplyCGOFromSnapshot}へ渡すのは入口ホップ一個の状態ではなく、式~\eqref{eq:cgo-state-sequence}の全状態列である。

\begin{algorithm}[H]
\caption{TLSレコードとCGO層を順に処理する保存通信の復号}
\label{alg:traffic-recovery}
\small
\begin{algorithmic}[1]
\Require $\mathcal{R}$、コアダンプ$M_t$、入口リンクを含む保存PCAP $P$
\Ensure TLSおよびCGOの検証結果を追加した$\mathcal{R}$
\ForAll{$R_c\in\mathcal{R}$かつ所有・経路・暗号コンテキストを検証済みで、全ホップがCGO}
  \State $T_c\gets\Call{RecoverEntryTLS}{M_t,c}$
  \If{$T_c$のTLS版、役割、暗号スイート、秘密長、回路帰属を検証できる}
    \State $\mathcal{D}_c\gets\Call{AuthenticateExactSequence}{P,T_c,T_c.\mathrm{write\_seq}}$
    \State $\mathcal{C}_c\gets\Call{EnumerateRelayCells}{\mathcal{D}_c,I_c.\mathrm{link\_circid}}$
    \If{$\mathcal{C}_c\neq\emptyset$}
      \State TLS AEAD認証とセル抽出の成功を記録; $U_c\gets(T_c.\mathrm{write\_seq},\mathcal{D}_c)$
      \State $\mathcal{A}_c\gets\emptyset$
      \ForAll{$q\in\mathcal{C}_c$をPCAP内の順に走査}
        \State $\widehat{\mathcal{S}}\gets\Call{Copy}{\mathcal{S}_{c,\mathrm{forward}}(t)}$
        \State $(h,z)\gets\Call{ApplyCGOFromSnapshot}{q,H_c,\widehat{\mathcal{S}}}$
        \If{$h$でCGO認証タグを認識}
          \State $\mathcal{A}_c\gets\mathcal{A}_c\cup\{(q,h,z)\}$
        \EndIf
      \EndFor
      \If{$|\mathcal{A}_c|=1$}
        \State CGO認証成功を記録; $Z_c\gets\operatorname{only}(\mathcal{A}_c)$
      \EndIf
    \EndIf
  \EndIf
\EndFor
\State \Return $\mathcal{R}$
\end{algorithmic}
\end{algorithm}

全候補でCGO認証タグの検証に失敗した場合、または複数候補で成功した場合、宛先ホップとリレーセル本体は一意に確定できない。構造検査、TLS認証、CGO認証は、相互に異なる状態と鍵に基づく検証条件である。これらの統計的独立性は仮定しない。平文復元が成立するには三つの検証がすべて成功する必要があり、いずれか一つの成功だけでは成立しない。

\subsection{検証項目と判定条件}

対象は、調査権限のある端末で動作するC Torクライアントである。
入力は、(1) ELFコアダンプ、(2) Build IDを記録し、コア内の実行ファイル写像との互換性を検証した実行ファイルとDWARFである。TLSおよびCGOの認証を検証する場合は、(3) 取得時点をまたぐ入口リンクのPCAPも必要となる。
解析は取得後に行い、動的デバッガや特別な識別子を対象プロセスへ埋め込まない。

復元成功は、単に候補が出たことではなく、次の検査を通過したことと定義した。

\begin{itemize}
  \item 回路がTor自身のorigin circuit所有リストに含まれる。
  \item list index、magic、ポインタ範囲、\code{next}/\code{prev}、閉路が整合する。
  \item lifecycleと暗号状態の初期化が整合する。
  \item PCAP復号ではTLS AEAD認証、Tor回路ID、CGOタグ認識を順に満たす。評価実験では、復元本文を別に取得した正解平文と比較する。
\end{itemize}

追加計測から得たControl Protocol記録、動的フックのセル、合成identityは評価専用の正解であり、解析入力から除外する。

各検査は表~\ref{tab:evidence-levels}の項目ごとに、成功、失敗、未評価のいずれかを記録する。後段の処理は、前提となる検査がすべて成功した場合だけ実行する。CGO認証タグの検証だけが成功しても、所有リストと\code{cpath}を検証できなければ、回路復元の成立条件を満たさない。

\begin{table}[!tbp]
\centering
\caption{解析結果として記録する検証項目}
\label{tab:evidence-levels}
\small
\begin{tabularx}{\linewidth}{L{34mm}Y Y}
\toprule
検証項目 & 成功条件 & 調査上言えること \\
\midrule
所有リストの完全列挙 & 所有リストの領域、要素数、容量、全要素を読み取れる & 対象ポインタが取得時点の所有リストに含まれる \\
回路構造と経路の検証 & magic、index、双方向リンク、閉路が整合 & 入口起点の経路構造を復元できる \\
暗号コンテキストの確認 & Tor内部のOPEN状態と方式別コンテキストの初期化が整合 & ホップ別の確立済み暗号方式を確認できる \\
TLSレコードのAEAD認証 & 入口SSLの秘密とコア内レコード番号で認証に成功 & PCAP内のTorセルを入口リンクへ対応付けられる \\
CGO認証タグの検証 & 成功候補が一件だけで、タグを除く493バイトを取得 & CGO認証に成功したホップとリレーセル本体を復元できる \\
\bottomrule
\end{tabularx}
\end{table}

所有リスト、回路構造、暗号コンテキストの検査にはコアダンプだけを用いる。TLSおよびCGOの認証には、コアダンプ取得時点をまたぐPCAPも必要である。後段の検証が成立しない場合でも、前段までに成立した結論は有効である。例えば、OpenSSLのメモリ配置が未対応でTLS秘密を抽出できない場合でも、所有リストと経路構造の検証結果は成立する。


取得したメモリを時点$t$のスナップショット$M_t$とする。$M_t$には、取得時点の所有構造、終了後の残存構造、未初期化領域、無関係な暗号コンテキストが混在し得る。本手法は、構造体の形や鍵のエントロピーを根拠に候補を集めるのではなく、Torが時点$t$に所有していたクライアント起点回路の集合$O_t$を最初に求める。

回路$c$を調査対象として受理する条件は、$c\in O_t$であり、かつ所有リスト上の位置と$c$自身が持つ所属情報が一致することである。この条件を満たさないメモリ断片は、後続のフィールドがもっともらしくても、取得時点の所有回路へ帰属できない。

この段階で立証するのは、回路が取得時点にTorの管理下にあったという事実である。通信が成功したこと、特定方式が使われたこと、セル本文を復号できることまでは主張しない。後続の各段階は、この所有証拠へ追加の条件を積み重ねる。


origin circuitの経路は循環リストである。循環リストだけを見れば、任意の要素を先頭にできるため、アドレスの大小や走査で最初に見つかった位置から入口を推定してはならない。本手法は、origin circuitが入口ホップとして保持する参照を$p_0$とし、$p_0,p_1,\ldots,p_{n-1}$を\code{next}関係で並べる。

並びを受理するには、各$p_i$で前後のリンクが相互に戻り、最後の要素から$p_0$へ閉じる必要がある。途中でNULL、別の閉路、読取り不能領域へ達した場合、確認できるのは所有リストへの掲載までであり、入口起点経路は確定できない。

この方法は「入口らしさ」を統計的に推定するものではない。Tor実装が持つ所有参照の意味を用いて向きを確定する。そのため、同じ三つのホップを得た場合でも、順序が逆または回転していれば成功とは数えない。


暗号状態について、次の三つを区別する。

\begin{description}[style=nextline,leftmargin=2.5em]
  \item[交渉意図]
  対向リレーの能力とハンドシェイク情報から、CGOを選択しようとしたことを示す。これは暗号コンテキストが完成したことを保証しない。

  \item[確立済み暗号方式]
  ホップが確立状態であり、方式識別子と必要な前方向・後方向の暗号コンテキストが整合することを示す。この条件が成立した場合だけCGOまたは\torone{}を確定する。

  \item[CGO累積状態の更新]
  CGOの累積状態$T'$が初期値から変化したことを示す。これはその方向で状態遷移が起きた証拠であるが、特定ストリームの通信、正確なセル数、ユーザ行為までは示さない。
\end{description}

この分離は構築中の回路に重要である。ゼロ初期化された列挙値が\torone{}と同じ値でも、暗号コンテキストが未初期化なら「\torone{}確立」とは判定しない。また、終了後に鍵断片が残っていても、所有集合$O_t$に回路がなければ取得時点の所有回路へ結び付けない。

\subsection{取得時点の暗号状態による通信復元}

CGO状態はセルごとに更新される。時点$t$のコアダンプが保持する対象回路・方向・ホップの状態を$S_{c,d,h}(t)$とする。この状態を適用できるのは、同じ状態を消費する次のセルである。したがって、PCAP内に同じ回路IDのセルがあるというだけでは不十分である。コアダンプ内のTLSレコード番号、入口リンクの回路ID、CGO認証タグを用いて、コアダンプ取得時点に対応するセルを検証する。

保存通信を解析する場合は、次の条件をすべて要求する。

\begin{enumerate}
  \item PCAPがコアダンプ取得時点$t$をまたいで入口リンクを記録している。
  \item コアダンプ内のクライアントTLS書込みレコード番号と、PCAPでAEAD認証できたクライアント側レコードの番号が一致する。
  \item そのレコード内で、対象入口回路IDを持つ全リレーセルを出現順に列挙する。
  \item 各候補について全ホップ状態列$\mathcal{S}_{c,\mathrm{forward}}(t)$を複製し、入口から順に各CGO層へ適用する。CGO認証タグの検証に成功した候補をすべて集め、成功候補が一件だけの場合に限り採用する。
\end{enumerate}

同じTLSレコードに別回路のセルが先に存在しても、対象回路のCGO状態は進まない。一方、対象回路のセルがコアダンプ取得とPCAPレコードの間に一つでも欠落すれば、それ以降のCGO状態との同期を検証できない。このため、欠損後のセルは復号対象にならない。

壁時計の時刻は補助情報であり、復号対象を決める主根拠ではない。コアダンプ内のTLS書込みレコード番号、入口リンクの回路ID、CGO認証という三つの証拠を用いる。対象回路の候補は認証済みTLSレコードからすべて列挙し、CGO認証に成功する候補が一件だけの場合に確定する。壁時計の時刻またはPCAP内の出現順だけでは、取得時点に対応するレコードとセルを確定できない。

\begin{figure}[!tbp]
\centering
\begin{tikzpicture}[
  x=1mm,y=1mm,>=Latex,
  icon/.style={minimum width=18mm,minimum height=15mm,
    inner sep=0pt,font=\fontsize{20}{20}\selectfont,text=TorPurple},
  label/.style={font=\small,align=center,text width=36mm,inner sep=0pt},
  check/.style={font=\footnotesize,align=center,fill=white,inner sep=1mm,text=black!75},
  arrow/.style={->,thick,TorPurple,shorten <=1mm,shorten >=1mm}
]
\node[font=\bfseries\small,text=TorPurple] at (-52,69) {クライアントメモリ};
\node[font=\bfseries\small,text=TorPurple] at (52,69) {保存PCAP};

\node[icon,text=black!65] (core) at (-35,52) {\faMemory};
\node[label,anchor=east] at (-46,52) {クライアントのコアダンプ\\実行ファイル・DWARF};
\node[icon,text=black!65] (pcap) at (35,52) {\faNetworkWired};
\node[label,anchor=west] at (46,52) {方向別TCPストリーム\\保存パケット};

\node[icon] (owner) at (-35,17) {\faRoute};
\node[label,anchor=east] at (-46,17) {Torが所有するクライアント起点回路\\回路ID・入口起点cpath};
\node[icon,text=black!65] (record) at (35,17) {\faStream};
\node[label,anchor=west] at (46,17) {対象TLSレコード\\方向・番号・暗号スイート};

\node[icon,text=black!80] (tls) at (-35,-18) {\faKey};
\node[label,anchor=east] at (-46,-18) {入口TLS 1.3状態\\通信秘密・IV・レコード番号};
\node[icon,text=black!80] (cell) at (35,-18) {\faServer};
\node[label,anchor=west] at (46,-18) {認証済みTorセル\\回路ID・セル順序};

\node[icon] (cgo) at (-35,-53) {\faLayerGroup};
\node[label,anchor=east] at (-46,-53) {ホップ別CGO状態\\$K,N,T'$・方向};
\node[icon,text=black] (plain) at (35,-53) {\faUnlock};
\node[label,anchor=west] at (46,-53) {CGO認証に成功した宛先ホップ\\リレー本文493バイト};

\draw[arrow] (core) -- node[right=2mm,check]{所有検証} (owner);
\draw[arrow] (owner) -- node[right=2mm,check]{入口TLS接続を特定} (tls);
\draw[arrow] (tls) -- node[right=2mm,check]{同じ回路} (cgo);

\draw[arrow] (pcap) -- node[left=2mm,check]{レコード抽出} (record);
\draw[arrow] (record) -- node[left=2mm,check]{順序照合} (cell);
\draw[arrow] (cell) -- node[left=2mm,check]{CGO層処理} (plain);

\draw[arrow] (tls) -- node[above=1mm,check]{TLS 1.3 AEAD認証復号} (cell);
\draw[arrow] (cgo) -- node[above=1mm,check]{CGOタグ検証} (plain);
\end{tikzpicture}
\caption{保存PCAPとクライアントメモリの対応関係。各段階が次段階の対象と検証条件を与える。}
\label{fig:chain}
\end{figure}


PCAPからリレー本文までには、相互に異なる鍵と状態を用いる二つの認証判定がある。第一はTLS~1.3のAEAD認証タグである。通信秘密、ノンス、レコード番号、TLSヘッダのいずれかが対応しなければ、レコードの認証は成立しない。第二はCGO認証タグである。入口から順に各ホップの状態を適用し、タグを認識したホップを宛先とする。両判定の統計的独立性は仮定しない。

TLS認証だけが成功した場合に確定できるのは、PCAP内のTorセルと入口リンクの対応までである。CGO認証タグの検証に成功したホップがなければ、宛先ホップとリレーセル本体は確定できない。状態ずれ、セル欠損、方向の誤り、未対応実装などが原因として考えられる。成功候補が一件だけの場合に限り、認証タグ16バイトを除く493バイトがリレーセル本体として確定する。

評価用平文との一致は、実験で正しさを測るための外部判定であり、候補選択には使用しない。実際の調査では正解平文がなくても、TLS AEADとCGO認証タグという二つの検証結果を得られる。

\enlargethispage{2\baselineskip}

本手法は、所有リストの完全列挙、回路構造、暗号コンテキスト、TLS AEAD認証、CGO認証タグの各検証結果を個別に返す。成功した検査については、使用した回路、方向、取得時点も記録する。

検証失敗から対象の不存在は導けない。所有リストを読み取れない場合、回路の不存在は判定できない。同様に、CGO認証タグが一致しない場合、通信の不存在は判定できない。各段階の結論は、必要な証拠がすべて得られた場合に限って確定する。証拠が不足する場合は当該段階を未評価とし、不足した証拠を記録する。

\section{実装}\label{sec:implementation}

評価実装は、GDBのPython APIでコアダンプを読み、DWARF型情報と対象バージョンで検証したOpenSSLのメモリ配置に基づいて内部状態を解釈する~\cite{gdbcorepython,dwarf5,openssl3022}。環境に依存するのは、型情報の解決方法とOpenSSL内部オフセットである。回路所有、入口順、コアダンプ取得時点との対応、TLSおよびCGO認証の判定条件は、これらの実装選択に依存しない。

\subsection{解析手順と実行ファイルの同定}

解析は表~\ref{tab:pipeline}の五段階に分けた。各段階は、前段の出力に含まれるアドレスまたは識別子だけを次段へ渡す。したがって、後段で暗号鍵らしい値が見つかっても、所有回路まで逆向きに結び付かない候補は採用しない。

\begin{table}[!tbp]
\centering
\caption{解析手順と各検証条件}
\label{tab:pipeline}
\small
\begin{tabularx}{\linewidth}{L{27mm}YY}
\toprule
処理 & 主な検査 & 後続処理へ渡す情報 \\
\midrule
型情報の解決 & Build IDを記録し、コア写像、DWARF、PIEロードバイアスを照合 & 型、フィールド位置、所有リスト位置 \\
回路構造の復元 & 所有リスト、magic、index、閉路を検証 & 回路ID、入口\code{cpath}、入口チャネル \\
暗号状態の抽出 & 暗号方式、初期化状態、OpenSSLのメモリ配置を検証 & TLS秘密、レコード番号、ホップ別CGO状態 \\
TLSレコードの復号 & TCP再構成、TLS AEAD認証、Torセル抽出 & 509バイトのリレーセルペイロード \\
CGO層の処理 & 入口から各層を処理し認証タグを検証 & CGO認証に成功したホップ、493バイトのリレーセル本体 \\
\bottomrule
\end{tabularx}
\end{table}

保存済みコアダンプの解析は読み取り専用であり、対象プロセス内の関数呼出しと値の書換えを伴わない。JSON形式の結果には、回路ごとの所有根拠、向き、構造検査、暗号初期化、エラー、完全性を含める。


候補実行ファイルのSHA-256とELF Build IDを記録する。Linuxコアは実行ファイルのBuild IDを常に直接保持するとは限らないため、\code{NT_FILE}対応表から主実行ファイルの写像を求め、コアに収録されたファイル由来バイトと候補ELFの同一offsetを比較する。\code{gcore}ではcleanな実行ページが省略され、再配置・可変データは変化するため、比較可能な110,592バイトのうち75\%以上の一致を互換性条件とした。正しいELFは84.51\%一致し、同版LTOと異版は8.71\%、9.58\%だった。この75\%閾値は本取得形式における対応ELFの受理条件であり、Build IDに代わる一般的な同一性証明ではない。

通常ビルドでは、\code{global_origin_circuit_list}の位置と各構造体のフィールド位置をDWARFから得る。このため、回路解析はTorバージョンや最適化条件ごとの固定オフセット表を必要としない。

LTOでは所有リスト変数のDWARF locationがoptimized outになる場合がある。この場合だけ、同一ELFのローカルシンボルと、実行時にも解決できる関数\code{circuit_get_global_origin_circuit_list}をanchorとしてPIE load biasを求める。ELF上のanchorアドレスを$A_{\mathrm{elf}}$、core内の実行時アドレスを$A_{\mathrm{run}}$とすると、

\begin{equation}
 b=A_{\mathrm{run}}(\mathrm{anchor})-A_{\mathrm{elf}}(\mathrm{anchor}),\qquad
 A_{\mathrm{run}}(s)=A_{\mathrm{elf}}(s)+b
\end{equation}

である。候補シンボルが0個または複数なら解析を停止する。外部から明示したload biasを使う場合も、対象ELF内でowner-list symbolが一意であることを要求する。誤バイナリを推測で補正しない。

\subsection{所有回路・経路・ストリームの復元}

解析器はDWARFから型、フィールド位置、enum、グローバル変数を解決する。
最初に\code{global_origin_circuit_list}を読み、全要素を列挙できたかを記録する。その後、各\code{origin_circuit_t}のmagicとリスト添字を確認する。
\code{smartlist}について、\code{num_used}が0以上、\code{capacity}以下であり、\code{capacity}が$10^6$以下であることを確認する。範囲外なら列挙しない。要素$i$を\code{origin_circuit_t*}へ変換した後、\code{base_.magic}が\code{ORIGIN_CIRCUIT_MAGIC}であることと、構造体内の\code{global_origin_circuit_list_idx}が$i$に一致することを検査する。入口起点閉路まで成立した要素だけを、構造検証済みの回路とする。

各回路では\code{origin_circuit_t.cpath}を入口ホップとする。
以後、最大32ホップまで\code{next}をたどり、各ノードについて次の不変条件を検査する。

\begin{center}
\code{hop.next.prev == hop},\quad
\code{hop.prev.next == hop},\quad
\code{hop.magic == CPATH_MAGIC}.
\end{center}

既訪問アドレスへ達したとき、そのアドレスが最初の\code{cpath}と同じ場合だけ閉路成立とする。NULL、読取り不能、別の既訪問ノード、32ホップ超過は構造エラーとして記録する。
入口は所有参照によって決まり、ホップのアドレス順には依存しない。

回路ごとに、Control Protocolと対応可能な\code{global_identifier}、入口リンクの\code{n_circ_id}、\code{n_chan}のチャネルIDと対向識別子を取得する。これらは探索条件ではなく、後続のTLSオブジェクトとPCAP内Torセルを同じ回路へ結ぶ結合キーになる。

所有リストを完全列挙でき、回路構造を検証でき、活動中の全ホップでCGO暗号コンテキストが初期化されている回路だけを、後続処理の対象回路一覧へ記録する。この一覧には、所有リスト添字、回路アドレス、二種類の回路ID、入口起点のホップアドレス列、ノンスと$T'$のハッシュ、元の構造解析結果のSHA-256を含める。TLS秘密抽出器とCGO秘密抽出器は同じ一覧を必須入力とし、所有リスト要素、回路アドレス、回路ID、ホップ数、ホップアドレスが一致しなければ処理を終了する。したがって、所有関係、経路構造、暗号コンテキストを検証していない候補が、TLSまたはCGOの処理だけで採用される経路はない。


クライアント起点回路が所有する\code{p_streams}から\code{edge_connection_t}の単方向リストをたどる。各要素について\code{on_circuit}が現在の回路アドレスと一致すること、入口接続のmagic、ストリームID、\code{cpath_layer}を確認する。SOCKS要求が存在する場合は、ハンドシェイク完了、宛先、ポートを取得する。

ストリームの帰属には、ヒープ全体のSOCKS構造体ではなく、Torが保持する\code{p_streams}所有辺を用いる。ユーザ名とパスワードは回路帰属に不要であるため収集対象から除外する。接続元端点は、調査目的上必要な場合に限って記録する。

\subsection{CGOと入口TLS状態の抽出}

各\code{crypt_path_t}から\code{relay_crypto_t}へ進み、暗号方式enumだけでなく初期化済みかを確認する。
\codeを\torone{}、\codeをCGOと解釈し、それ以外は未知値とする。構築中のゼロ初期化メモリでは、列挙値のゼロが\torone{}と同じ値に見える場合がある。そのため、回路とホップがTor内部のOPEN状態であり、対応する暗号コンテキストが初期化済みの場合だけ方式を確定する。\torone{}では前方向・後方向の暗号器とダイジェストの4ポインタ、CGOでは前方向・後方向の\code{cgo_crypt_t*}がすべて非NULLであることを要求する。\code{extend_info.enable_cgo}が得られる場合は、交渉意図と確立済みの\code{kind}の一致も記録する。

CGOでは前方向と後方向を別々に処理し、表~\ref{tab:cgo-components}の96バイトを一方向の現行状態として取得する。Tor構造体が直接保持するPOLYVAL鍵、ノンス、$T'$はDWARFで解決する。AES鍵だけは\code{EVP_CIPHER_CTX}からOpenSSLプロバイダーのコンテキストをたどる。

\begin{table}[!tbp]
\centering
\caption{一方向のCGO現行状態として取得する成分}
\label{tab:cgo-components}
\small
\begin{tabularx}{\linewidth}{YcY}
\toprule
成分 & 長さ & 取得元と検査 \\
\midrule
$J$用AES-128鍵 & 16 B & プロバイダーのAES復号鍵スケジュール \\
$J$用POLYVAL鍵 & 16 B & \code{uiv.j.ku.key.h} \\
$S$用AES-128鍵 & 16 B & プロバイダーのAES暗号鍵スケジュール \\
$S$用POLYVAL鍵 & 16 B & \code{uiv.s.b.h} \\
ノンス$N$ & 16 B & \code{cgo_crypt_t.nonce} \\
累積状態 $T'$ & 16 B & \code{cgo_crypt_t.tprime} \\
\bottomrule
\end{tabularx}
\end{table}

OpenSSL~3の内部には公開ABIがない。暗号状態の抽出は、対象ライブラリが検証済みの\code{openssl-3.0.22-debian-arm64}プロファイルと一致する場合に限定する。このプロファイルでは\code{EVP_CIPHER_CTX.algctx}をオフセット168で読み、プロバイダーのAES鍵スケジュールを\code{algctx+192}から取得する。AESラウンド数は鍵スケジュール内のオフセット240で読み、AES-128を示す10でなければ拒否する。$J$の復号鍵スケジュールでは最終ラウンドの位置、$S$の暗号鍵スケジュールでは最初のラウンドの位置から16バイトを取得した。この対応は同じビルドの単体試験が保持する現行鍵と照合して検証した。

\code{algctx == NULL}、AESラウンド数の不一致、各成分が16バイトでない場合、読み取り不能ページがある場合は、その方向を復号に使わない。標準出力には状態全体と$T'$のSHA-256だけを記録する。生の96バイトを証拠として保存する場合は、ファイル権限を0600とする。
また、$T'\neq 0$は、その方向で状態遷移が起きた証拠として扱う。特定ユーザの通信やセル数を示す証拠とは扱わない。


所有確認済みorigin circuitから、次の型付きポインタ列をたどる。

\begin{center}
\small
\code{origin_circuit_t} $\rightarrow$ \code{base_.n_chan} $\rightarrow$
\code{channel_tls_t.conn} $\rightarrow$ \code{or_connection_t.tls}
$\rightarrow$ \code{tor_tls_t.ssl} $\rightarrow$ \code{SSL}
\end{center}

途中のポインタがNULLなら、対象回路と入口TLS状態の対応は検証不能となる。\code{SSL.version}が\code{0x0304}、\code{SSL.server}と\code{tor_tls_t.isServer}がともにクライアント側を示すことを確認する。検証済みOpenSSLプロファイルからクライアント／サーバのアプリケーション通信秘密、読込み／書込みIV、読込み／書込みレコード番号、クライアント乱数を読む。

暗号スイートは通信秘密の長さから推測しない。対象OpenSSLソースとディストリビューション設定から\code{SSL.session}、\code{SSL_SESSION.cipher}、\code{SSL_CIPHER.id}の位置を導出し、コアダンプ内の識別子を直接読む。本評価では全対象が\code{0x1302}、すなわち\code{TLS_AES_256_GCM_SHA384}であった。識別子、ハッシュ、秘密長が対応しない場合、または実装していない\code{TLS_CHACHA20_POLY1305_SHA256}の場合は復号へ進まない。

\begin{table}[!tbp]
\centering
\caption{OpenSSL 3.0.22 Debian arm64用TLSプロファイルの主要オフセット}
\label{tab:tls-profile}
\small
\begin{tabular}{lr@{\qquad}lr}
\toprule
フィールド & オフセット & フィールド & オフセット \\
\midrule
\code{SSL.version} & 0 & \code{SSL.server} & 56 \\
client random & 352 & \code{SSL.session} & 2328 \\
client app secret & 1860 & server app secret & 1924 \\
read IV & 2128 & write IV & 2176 \\
read sequence & 7416 & write sequence & 7424 \\
\code{SSL_SESSION.cipher} & 760 & \code{SSL_CIPHER.id} & 24 \\
\bottomrule
\end{tabular}
\end{table}

通信秘密領域64バイトの末尾がゼロであることを確認し、非ゼロ部分を32または48バイトとして採用する。クライアント／サーバの長さが異なる場合と全ゼロの場合は拒否する。暗号スイート識別子から決まるハッシュと秘密長が実データに一致する場合だけ採用する。レコード番号はOpenSSLレコード層の8バイトbig-endian値として読む。

\subsection{TLSレコードとCGO層の復号条件}

コアダンプ取得後の保存通信の復号は、図~\ref{fig:chain}の順に行う。PCAPからEthernet/IPv4/TCPペイロードを取り出し、送信元・宛先IPとポートの四つ組ごとにTCP番号順で再構成する。同じ番号の再送は最初のペイロードを一つだけ残す。再構成したストリームから5バイトのTLSヘッダとレコード本体を切り出す。

アプリケーション通信秘密から、式~\eqref{eq:tls-key}に従ってAES-256-GCMの鍵とノンスを導出する。

外側のTLSヘッダを関連データとし、コアダンプが保持するクライアント側書込みレコード番号をノンスへ用いてGCM認証タグを検証する。0--64の総当たりは行わない。正しいTLS秘密でも別のレコード番号を与えれば、その番号に対応する後続レコードは暗号学的に正しく認証され得る。このため、AEAD認証だけで取得時点を判断せず、コアダンプ内番号と取得順序の一致も必須条件とする。

TLSInnerPlaintextの末尾パディングと内部コンテンツタイプを除き、アプリケーションデータからTorセルを探す。通信路上のセル形式は、4バイトの入口リンク回路ID、1バイトのコマンド、509バイトのペイロードである。コマンドは\code{RELAY=3}または\code{RELAY_EARLY=9}に限定した。同じTLSレコード内に対象回路のセルが複数ある場合、Control Protocol由来の正解回路IDや評価用正解セルで一件へ絞らず、出現順に全候補を列挙する。


抽出した各ホップの前方向状態80バイトと$T'$ 16バイトから、Tor実装と同じ\code{cgo_crypt_t}を解析用プロセス内に再構成する。各セル候補について取得時点状態の複製を作り、509バイトをセルペイロードへ設定し、入口ホップから順に\code{cgo_crypt_relay_forward}相当の処理を適用する。全候補を処理した後、認証タグを認識した候補が厳密に一件である場合だけ、その候補とホップを採用する。認証候補が0件または複数件なら状態不一致または曖昧として平文を出力しない。

CGO認証タグである先頭16バイトを除いた493バイトをリレーメッセージ本体として出力する。復号用の試験入力には、コアダンプ由来のCGO状態とPCAP由来の暗号文だけを格納し、評価用の正解ホップと正解平文を含めない。復号処理が終了した後、別プロセスの評価器だけが、動的フック由来の正解セルとCGO認証に成功したホップ、493バイトを比較する。この物理的分離により、正解情報を回路選択、TLS復号、試験入力、CGO処理へ混入させない。

解析結果には秘密値そのものを含めず、方式、初期化状態、更新有無、検証結果、秘密値のSHA-256を記録する。リレー本文は標準出力から除外し、証拠内容の取得が調査目的に含まれる場合だけ保存する。


解析結果が0件でも、回路が存在しなかったとは限らない。所有リストの\code{smartlist}を含むページが欠損すれば、解析器は候補を得られないからである。
そこで、必要なELFセグメント、所有リスト、各回路、各暗号経路、OpenSSL領域について、読み取りの成否と列挙の完了状態を出力する。
所有リストを全要素列挙できた場合は「列挙完了」、一部で読取りエラーが生じた場合は「部分復元」、所有リスト自体を解決できない場合は「取得不能」とする。\code{num_used=0}の場合に言えるのは、取得時点のリストが空であったことだけである。必須ページが欠けた場合は「0回路」ではなく「復元不完全」と報告する。

個別回路では、構造検査と暗号検査を分離する。所有回路として列挙できても\code{cpath}閉路が壊れていれば、回路メタデータは残し、暗号状態を確定しない。CGO構造体を読めてもOpenSSLプロファイル検査に失敗すれば、方式はCGO、秘密抽出は失敗として記録する。段階別の結果により、回路不存在、経路構造の欠損、暗号状態の取得失敗を区別できる。

標準JSONにはアドレス、方式、状態、検査結果、秘密値のハッシュだけを保存する。TLS秘密、CGO状態、復号済みTLSレコード、リレー本文を証拠として保存する場合は、生成ファイルのアクセス権を0600とする。公開用集計には生の秘密を含めない。

\section{実験設計}\label{sec:experiment}

\subsection{隔離環境と対象}

第三者通信を含まないChutney隔離網を構築し、制御下のクライアントとリレーだけを使用した。使用したChutneyとC Torのコミットを固定した~\cite{chutneysource,ctorsource04916}。
主要な評価環境はLinux arm64上のC Tor 0.4.9.16、OpenSSL 3.0.22である。
一般化評価ではC Tor 0.4.9.3-alpha、0.4.9.9、0.4.9.16、\code{-O0}/\code{-O2}/LTOを用いた。
x86\_64 ABIは、QEMUユーザーモードで取得したゲストのコアダンプで評価した。これはネイティブx86\_64評価ではない。
実環境への適用可能性を確認する評価では、Tor Project公式リポジトリがDebian 12 arm64向けに配布するTor 0.4.9.17と\code{tor-dbgsym}を用いた~\cite{torpackage04917}。実行ファイルはストリップ済みであり、GNU Build IDが一致する外部デバッグシンボルを型情報として使用した。

0.4.9.3-alphaと0.4.9.9には、通常の自動構築回路で最終ホップ以外のCGO交渉を妨げる既知不具合がある~\cite{torrelease04912}。バージョン・最適化差の評価では、明示経路を\code{EXTENDCIRCUIT purpose=controller}で構築して全ホップCGOを成立させた。この実験は解析器の評価であり、各版におけるCGOの自然な利用割合を測るものではない。この差を確認するため、Torが自動構築した回路の方式分布も別に集計した。

実験用ストリームの宛先、送信元、識別子は合成値とした。SOCKS認証情報は解析結果へ出力しない。メモリ、PCAP、正解記録にはSHA-256を付与し、実行環境とソースコードのコミット識別子を記録した。

通信を一時停止する実験の時系列を図~\ref{fig:acquisition-timeline}に示す。PCAPを先に開始し、ストリームを一時停止してからTorを\code{SIGSTOP}し、コアダンプを取得する。Torを再開した後にストリームを再開し、入口リンクへ出る次のTLSレコードを保存する。動的フックは評価用セルを別ファイルへ保存するだけであり、復号処理へは渡さない。

これとは別に、停止時点をセルやTLSレコードの境界に合わせない実験を行った。PCAPと継続ストリームを開始し、ストリームを一時停止せず、動的セルフックも使わず、0.11--1.61秒の異なる遅延後に\code{gcore}を実行した。\code{gcore}はアドレス空間の複写中に対象プロセスを停止するため、この停止は6試行すべてに含まれる。一方、停止時点の選択にはセル処理やTLSレコード境界を用いなかった。2ホップと3ホップを各3独立起動とし、復元処理は正解セルを受け取らない。CGO認証で候補を一意化した後だけ、制御下ストリームへ繰り返し挿入した評価用マーカーを確認した。

\begin{figure}[!tbp]
\centering
\begin{tikzpicture}[x=18mm,y=1mm,>=Latex,
  event/.style={font=\footnotesize,align=center,text width=20mm},
  arrow/.style={->,thick,TorPurple}]
\draw[arrow] (0,0) -- (7.4,0);
\foreach \x/\top/\bottom in {
  0/{PCAP開始}/{継続取得},
  1/{ストリーム停止}/{送信抑止},
  2/{Tor停止}/{\code{SIGSTOP}},
  3/{コア取得}/{状態$M_t$},
  4/{Tor再開}/{\code{SIGCONT}},
  5/{送信再開}/{次セル生成},
  6/{TLS記録}/{PCAPへ保存},
  7/{評価記録}/{別系統の正解}}
  {
    \fill[TorPurple] (\x,0) circle (1.2pt);
    \node[event,above=2mm] at (\x,0) {\top};
    \node[event,below=2mm,text=black!65] at (\x,0) {\bottom};
  }
\end{tikzpicture}
\caption{通信を一時停止する実験の時系列。評価記録は復号終了後の比較にだけ用いる。}
\label{fig:acquisition-timeline}
\end{figure}

表~\ref{tab:experiment-groups}に実験群と独立単位をまとめる。実験群ごとに母集団と独立単位が異なるため、結果は群別に集計した。

\begin{table}[!tbp]
\centering
\caption{実験群、目的、独立単位}
\label{tab:experiment-groups}
\small
\begin{tabularx}{\linewidth}{Y L{25mm} L{31mm}}
\toprule
目的と条件 & 独立単位 & 規模 \\
\midrule
0.4.9系・最適化差の構造復元 & Torプロセス & 12起動、144回路 \\
ライフサイクルと残存構造 & 回路／取得時点 & 確立12、構築中8、終了6 \\
TLSレコードとCGO層の連続復号 & Torプロセス & 10起動、各7候補 \\
端点間の時間差と直前セル & Torプロセス & 3起動、36ホップ \\
実行ファイル・欠損・誤状態の負の対照 & 条件 & ELF 3種、欠損2、誤状態7,168 \\
既知再同期攻撃の実装境界 & Torプロセス & 3起動 \\
ストリームを停止せず、セル・TLS境界を観測しない固定遅延取得 & Torプロセス & 6起動、2/3ホップ \\
CGO/\torone{}混在回路の後続処理への入力制限 & 回路 & 19回路、うち混在10 \\
公式Debianパッケージ版への適用 & Torプロセス & 構造復元3起動・36回路、通信復元6起動 \\
\bottomrule
\end{tabularx}
\end{table}

\subsection{評価項目}

評価は次の六群に分けた。

\begin{enumerate}
  \item 所有回路、入口起点パス、CGO/\torone{}混在の正確性。
  \item バージョン、最適化、LTO、ABIの差に対する再現性。
  \item 確立／構築中／終了後、欠損ページ、誤バイナリに対する誤受理抑制。
  \item ストリーム、回路、経路、CGO状態、次セルの結合。
  \item コアダンプとPCAPだけを入力とするTLSレコードおよびCGO層の復号。
  \item 取得時点のずれ、直前セル、誤鍵、既知の後方向再同期攻撃という境界。
\end{enumerate}

回路・ホップは同じプロセス内で相関する。したがって、反復成功率はTorプロセスの状態を毎回初期化した起動を単位とする。二側95\% Wilsonスコア区間は、各起動を交換可能なBernoulli試行とみなした記述的な区間として示す。ホップ数や取得遅延を変えた試行は同一分布からの無作為標本ではないため、この区間から一般環境の成功率を推定しない。決定論的に生成した誤状態には信頼区間を付けず、観測件数だけを報告する。

\section{結果}\label{sec:results}

\subsection{検証条件を満たす回路・暗号状態を復元した}

初期の実Tor評価では、0.4.9.16単独環境の20/20回路、0.4.9.16と0.4.8.25を混在させた環境の19/19回路がControl Protocolの回路IDと一致した。入口起点の全経路も39/39で一致した。
混在環境の51ホップは、CGO 40、\torone{} 11に分類された。0.4.8.25リレーを含むホップはすべて\torone{}であった。
この混在コアを提案手法で再解析すると、所有関係、経路構造、暗号コンテキストを検証できた活動中回路19件のうち、全ホップCGOは8件、CGO/\torone{}混在は10件、\torone{}のみは1件だった。TLS・CGO通信復元の成立条件を満たした回路は、全ホップCGOの8件と厳密に一致した。混在10回路と\torone{}のみの1回路は、全ホップCGO条件を満たさなかった。

より広い評価では、四つのarm64条件を各3回起動し、制御下144/144回路と432/432ホップを復元した（表~\ref{tab:generalization}）。12プロセスすべてが完全成功し、プロセス単位の95\% Wilson区間は75.75--100\%である。

\begin{table}[!tbp]
\centering
\begin{threeparttable}
\caption{バージョン・最適化差に対する復元結果}
\label{tab:generalization}
\small
\begin{tabularx}{\linewidth}{Yrrr}
\toprule
条件 & 完全成功 & 回路 & CGOホップ \\
\midrule
0.4.9.3-alpha, \code{-O0} & 3/3 & 36/36 & 108/108 \\
0.4.9.9, \code{-O0}       & 3/3 & 36/36 & 108/108 \\
0.4.9.16, \code{-O0}      & 3/3 & 36/36 & 108/108 \\
0.4.9.16, \code{-O2}      & 3/3 & 36/36 & 108/108 \\
\midrule
合計 & 12/12 & 144/144 & 432/432 \\
\bottomrule
\end{tabularx}
\begin{tablenotes}\footnotesize
\item 回路・ホップは同一プロセス内で相関するため、独立成功率の単位はプロセスとした。
\end{tablenotes}
\end{threeparttable}
\end{table}

\subsection{公式Tor 0.4.9.17への適用}

自己ビルドではない配布パッケージへの適用結果を表~\ref{tab:packaged}に示す。Tor Project公式リポジトリのDebian 12 arm64向けTor 0.4.9.17はストリップ済みであり、実行ファイルのGNU Build IDは\code{59de31af...a9f}だった。同じ公式リポジトリの\code{tor-dbgsym}に含まれる外部デバッグファイルは、このBuild IDと一致した。3回の独立起動で36/36回路の所有関係と入口起点経路、108/108ホップのCGO確立状態を復元した。プロセス単位3/3成功の二側95\% Wilsonスコア区間は43.85--100\%であり、少数条件の成立確認であって配布環境一般の成功率推定ではない。

\begin{table}[!tbp]
\centering
\caption{Tor Project公式Debianパッケージ版の評価}
\label{tab:packaged}
\small
\begin{tabularx}{\linewidth}{L{29mm}L{18mm}YY}
\toprule
評価 & 独立起動 & 構造・認証結果 & 復元内容 \\
\midrule
構造復元 & 3/3 & 回路36/36、CGOホップ108/108 & 所有関係と入口順 \\
保存通信の再生 & 6/6 & 一意なCGO認証6/6 & DATA 5件、SENDME 1件 \\
\bottomrule
\end{tabularx}
\end{table}

外部デバッグシンボルを与えない条件では型情報を解決できず、回路復元は成立しなかった。異なるBuild IDのデバッグファイルを配置した条件でもGDBが型情報を読み込まず、回路復元は成立しなかった。実行ファイルとコアダンプの対応検査では、正しい配布実行ファイルは比較可能な368,640バイトの95.57\%が一致して通過した。別ビルドのTor 0.4.9.16は39.62\%だけが一致し、互換性条件を満たさなかった。したがって、配布実行ファイルだけでは本手法を適用できないが、対応する外部デバッグシンボルを保全できれば、所有リスト、回路構造、暗号コンテキストを自己ビルド時と同じ条件で検証できる。

同じ配布実行ファイルとデバッグシンボルを用い、隔離網とTorプロセスを毎回作り直した別の6起動で通信復元を評価した。通信は一時停止せず、ストリーム開始後0.13、0.31、0.57、0.19、0.43、0.71秒の固定遅延後にコアダンプ取得を開始した。セル処理やTLSレコード境界は観測していない。2ホップと3ホップを各3回試し、コアダンプ内のTLS書込みレコード番号13、16、19、25、29、34に対応するレコードを各1件認証した。6/6起動でCGO認証に成功した候補は一件だけだった。2ホップではホップ1、3ホップではホップ2が認識された。所有関係の確認からTLS AEAD認証、一意なCGO認証までの全検査は6/6で成功した。交換可能なBernoulli試行と仮定した記述的な95\% Wilsonスコア区間は60.97--100\%である。

CGO認証済みの493バイトのうち、評価用マーカーを含むリレーDATAセル本体は5/6だった。2ホップでは2/3、3ホップでは3/3である。残る1起動で復元したリレーコマンドはSENDMEであり、CGO認証には成功したが評価用マーカーを含まなかった。評価用マーカーを含む本体5/6の記述的な95\% Wilsonスコア区間は43.65--96.99\%である。CGO認証成功は、調査者が意図した通信内容の取得を保証しない。

再生器はTor Projectの公式0.4.9.17ソースアーカイブ（SHA-256 \code{ff372c41...abe2}）から構築した。製品側ソースは変更せず、保存済み状態を読み込む試験用インターフェースだけを\code{test_crypto_cgo.c}へ加えた。再生器のSHA-256は\code{e71129d0...0acbd}、GNU Build IDは\code{cfebbca1...eebc}である。保存済み6起動の全36候補を再処理し、従来の0.4.9.16ベース再生結果と、CGO認証ホップ、復元バイト列、正解情報不使用フラグが36/36で一致した。

TLSの負の対照では、通信秘密を1ビット変更した条件と別起動のTLS状態を与えた条件で認証レコードは0件だった。同じ接続秘密で書込み番号を1増やすと後続の正しいレコードは認証された。従って、AEAD認証だけで取得直後のレコードとは決めず、コアダンプ内番号と取得順序の一致を必要条件とする。

一方、明示経路でない自動生成回路では、0.4.9.3-alphaの3起動でCGO 133ホップ、\torone{} 91ホップ、0.4.9.9でCGO 132、\torone{} 93だった。修正後の0.4.9.16ではCGO 438、\torone{} 12である。この分布差は、0.4.9.12で修正された「最終ホップ以外でCGOを交渉しない」条件と整合する~\cite{torrelease04912}。ただし自動生成回路はホップ選択と長さを固定していないため、この値をネットワーク全体の普及率とは解釈しない。表~\ref{tab:generalization}の432/432は、CGO暗号コンテキストの初期化が確認できた既知回路を、各版の解析器が復元できるかという条件付き評価である。

LTOビルドでは、DWARF上の所有リスト位置をGDBから直接取得できなかった。ELFローカルシンボルとPIEロードバイアスを用いる代替処理により12/12回路、36/36 CGOホップを復元した。
QEMU上のx86\_64ゲストコアでも12/12回路、36/36ホップが一致した。


確立状態12/12、構築中8/8、正常終了後に所有リストから不在6/6を分類した。
構築中の16ホップは列挙値だけを見ると\torone{}であったが、暗号状態は未初期化だった。このため方式は判定不能であり、\torone{}確立の条件を満たさなかった。

完全なELFコア内の4~KiBページをゼロ化した8条件では、偽の構造的有効候補は0件だった。この評価は真のページ欠損ではなく、ゼロ化した擬似欠損である。所有リストの\code{smartlist}をゼロ化した条件では、読取りの成否と列挙完了を検査しなければ、エラーなしの0回路と解釈される。しかし、この条件では所有リストの全要素を確認できないため、回路の不存在は判定できない。列挙完了を検査した結果は、所有リスト取得不能だった。

さらに、所有リストの\code{smartlist}を含むヒープ\code{PT_LOAD}と、所有リスト変数を含む実行ファイルのデータ\code{PT_LOAD}をそれぞれコアのプログラムヘッダから除外し、実際に読取り不能とした2条件を評価した。両条件とも所有リストは取得不能となり、回路構造を検証できなかった。

実行ファイル照合では、対応ELF、同版LTO、異版の3条件を用いた。対応ELFは写像バイト84.51\%一致で通過し、同版LTOと異版は8.71\%、9.58\%で互換性条件を満たさなかった。後二条件では型解決へ進まず、回路数は判定不能だった。

約731~MiBの完全なコアダンプに対する解析時間中央値は、制御回路0/12/48/95で0.142/0.164/0.223/0.306秒だった。最大RSSは57,404--60,224~KiBである。取得による停止時間は1.33--1.47秒、ハッシュ計算は0.92--1.04秒だった。

\subsection{コアダンプ取得後の通信を復元した}

クライアントのコアダンプから、制御下12/12ストリームの宛先、Control回路ID、入口起点経路を結合し、偽陽性・偽陰性はいずれも0だった。
同じ所有回路から抽出した現行CGO状態へ、コアダンプ取得後に生成された次のリレーセルを入力した。独立5起動すべてでホップ2を認識し、493/493バイトの本文が評価用平文と一致した。

負対照をCGO状態全体へ拡張した。$J$側AES鍵、$J$側POLYVAL鍵、$S$側AES鍵、$S$側POLYVAL鍵、ノンス、$T'$、ホップ順の7分類について各1,024状態、合計7,168状態を同じ実セルへ適用し、誤認識は0件だった。これらは決定論的対照であり、全誤状態からの無作為標本ではないため、偽陽性率の信頼区間へ一般化しない。
継続通信中の任意時点取得では、3プロセスから各10スナップショットを取り、CGO状態抽出は30/30で成功した。ただし同一プロセス内のスナップショットは相関する。同一プロセス内の10スナップショットを一群としてプロセス単位に集約すると、成功は3/3プロセスであり、95\%区間は43.85--100\%である。


復号入力は、取得時点が対応するクライアントコアダンプ、検証済み実行ファイルと型情報、入口リンクPCAPだけである。Control Protocol由来の正解回路IDと動的フック由来の正解セルは復号入力に含めず、候補確定後の照合だけに用いた。

保存済みの独立10取得物を、所有関係、経路構造、暗号コンテキストの検証を通信復元の必須条件として再解析した。10/10起動で所有リストを完全列挙し、三つの検査すべてに成功した。後続のTLSおよびCGO検証は、これらの条件を満たした全ホップCGO回路に対して実行した。

迂回可能性の負の対照として、暗号コンテキストを検証できた回路が存在しない条件、通信復元に必要な回路情報が欠けた条件、回路アドレスを改変した条件を与えた。4条件すべてで検証は失敗し、秘密値の復元は0件だった。

各取得物ではOpenSSL状態から暗号スイート\code{TLS_AES_256_GCM_SHA384}を直接取得し、コアダンプ内の書込みレコード番号7に対応するTLS~1.3レコードを1件認証復号した。対象入口回路IDを持つ7セルを、Control Protocol由来の選択子を使わずに列挙した。

全候補へ同じ取得時点のCGO状態の複製を適用した結果、各試行でCGO認証タグの検証に成功した候補は厳密に1件であり、その宛先はホップ2だった。候補確定後に正解情報と照合し、認証候補が正解候補と一致すること、PCAP由来509バイトと動的正解、認証タグを除く493バイトと正解平文が一致することを確認した。所有リストの完全列挙から正解平文との事後照合まで、表~\ref{tab:full}の各検査は10/10で成功した。

\begin{table}[!tbp]
\centering
\caption{TLSレコードとCGO層を順に処理した検証結果}
\label{tab:full}
\small
\begin{tabularx}{\linewidth}{YcL{31mm}}
\toprule
検証項目 & 結果 & 成功条件 \\
\midrule
所有リストの完全列挙 & 10/10 & 全要素を読取り \\
回路構造と入口起点経路 & 10/10 & magic、index、双方向リンク、閉路が整合 \\
CGO暗号コンテキスト & 10/10 & 活動中の全ホップで初期化と構造が整合 \\
TLS AEAD認証とセル抽出 & 10/10 & レコード番号7で各1レコード、各7セル \\
一意なCGO認証 & 10/10 & 認証候補が各1件、CGO認証成功ホップ2 \\
評価用正解との事後一致 & 10/10 & 候補、ホップ、本文493/493バイトが一致 \\
\bottomrule
\end{tabularx}
\end{table}

独立単位をTorプロセスとすると、固定条件下の反復成功率10/10の二側95\% Wilson区間は72.25--100\%である。
ストリームを止めず、セル・TLS境界を観測しない固定遅延で取得した追加実験でも、6/6起動で所有関係の確認から一意なCGO認証までの全検査に成功した。2ホップ3起動ではCGO認証成功ホップ1、3ホップ3起動ではホップ2となった。コアダンプから得たTLS書込みレコード番号は17、23、29、37、53、68と全試行で異なり、各番号でTLS~1.3レコード1件を認証した。各試行で得た7セル候補のうちCGO認証に成功したのは一件だけであり、その後に確認した評価用マーカーも一致した。6/6の記述的な95\% Wilsonスコア区間は60.97--100\%である。

この追加実験ではControl Protocol回路ID、動的フックのセル、正解ホップ、正解平文を復号入力に含めなかった。したがって、固定レコード番号7、固定ホップ2、ストリーム停止によってだけ10/10が成立した可能性は棄却できた。ただし1ホップの一般SOCKSストリームはTorが当該回路を通常の出口回路として使用せず、ホップ0を宛先とする通信復号は評価できなかった。誤鍵、別起動秘密、暗号コンテキストの検証失敗、検証済み回路情報の改変、混在回路の各条件では、リレーセル本体の復元は0件だった。一方、この6試行は全件成功したため、固定遅延取得の不成立例は観測していない。

TLS秘密を含まないPCAPだけを与えた対照では、118パケット、TCPペイロード78セグメントを再構成しても、既知セル、平文、合成識別子のいずれも回収できなかった。

TLS側の負の対照では、クライアント側通信秘密の1ビット反転と、別起動のTLS秘密による認証レコードはいずれも0件だった。一方、同じ接続秘密のまま書込みレコード番号を7から8へ変えると、後続の番号8のレコード自体は正しく認証された。AEAD認証はレコードの暗号学的正しさを示すが、それが取得直後のレコードであることまでは示さない。この結果から、TLSレコードの採用には、コアダンプ内レコード番号が改変されていないことと、図~\ref{fig:acquisition-timeline}の取得順序を必要条件とする。

この結果は、メモリだけ、またはPCAPだけで復号したという主張ではない。時点が一致する両方の証拠を、回路所有関係により結合した結果である。

\subsection{時間差は構造対応より先に暗号状態を失わせた}

100~MBのストリームを12回路で転送中に、クライアントと3リレーを順次取得した。
実測2.297、2.663、3.168秒の取得差でも、リンク列は36/36ホップで対応した。一方、ノンスと$T'$の完全一致は0/36だった。
回路構造は秒単位で持続しても、暗号状態はセルごとに進むためである。

また、セル処理直後の現行状態を直前の前方向セルへ適用した3回は、すべてCGO認証成功ホップが$-1$となり、平文も一致しなかった。この結果は図~\ref{fig:time}の時間境界を実Torで確認したものである。

\begin{table}[!tbp]
\centering
\caption{成立した復元と成立しなかった復元}
\label{tab:boundary}
\small
\begin{tabularx}{\linewidth}{YcY}
\toprule
問い & 結果 & 解釈 \\
\midrule
所有回路・入口起点経路 & 144/144 & 構造上の対応を再現 \\
コアダンプ取得後の次セル & 5/5 & 現行CGO状態が利用可能 \\
コアダンプ＋PCAP & 10/10 & TLSレコードからリレー本文まで復号 \\
ストリーム非停止・固定遅延取得 & 6/6 & TLS番号6種、宛先ホップ1/2で一意なCGO認証 \\
公式配布版・一意なCGO認証 & 6/6 & マーカーを含むDATAセル本体5/6、マーカーを含まない本体1/6 \\
混在回路の通信復元条件への適合 & 0/10 & 全ホップCGO条件を満たさない \\
現行状態＋直前セル & 0/3 & 過去状態を一般には逆算不可 \\
順次取得した両端状態 & 0/36 & 瞬時照合には同時停止が必要 \\
誤CGO状態の宛先誤認 & 0/7168 & 7分類の決定論的負対照 \\
\bottomrule
\end{tabularx}
\end{table}

\section{抽出状態が与える攻撃能力の境界}\label{sec:security}

メモリから暗号状態を抽出できても、それだけで遠隔攻撃や匿名性破壊が可能になるとは限らない。この脅威境界を確認するため、Pillwaxが示した後方向CGOに対するIPR+再同期識別攻撃をC Tor 0.4.9.16で再現し、実際の3ホップ回路でどこまで進むかを調べた~\cite{pillwax2026improving}。これは提案手法とは別の脆弱性主張ではなく、抽出状態から外挿してはならない能力を測る境界実験である。

単一中間層では、改変セルの左側タグを保持したまま右側を1ビット変更すると、そのセルは認識されないが、クライアント側とリレー側の対象層は再同期した。次の正常セルは509/509バイトが一致した。
したがって、形式的な再同期性は実装にも存在する。

しかし、3ホップ処理では外側の対象層だけが再同期し、改変セルを処理した内側2層の$T'$はリレー側と不一致になった。回路を仮に維持しても、次の正常な出口セルは認識されなかった。同じ2セル列はArti 0.48.0でも\code{BadCellAuth}として拒否された~\cite{artisource048}。
さらに、隔離Chutney網で入口リレーがセルを1回だけ改変した3回の独立起動では、3/3でクライアントが対象回路を閉じた。95\% Wilson区間は43.85--100\%であり、小標本である。

以上から、既知の形式的攻撃は実装で再現できるが、本評価では鍵なし復号、セル偽造、匿名性破壊へ発展しなかった。全層処理による内側状態の不一致と、最初の認証失敗で回路を閉じる処理が、2回目の問い合わせを妨げた。本研究はこれを新規脆弱性として主張しない。

\section{考察}\label{sec:discussion}

\subsection{実務上の意味と能力境界}

本手法が直接答えられる問いは具体的である。
取得時点にTorが所有していたクライアント起点回路、そのControl回路ID、入口からのリレー順、ホップごとのCGO/\torone{}、方向別CGO状態の初期化と更新有無を復元できる。
ストリーム構造が残っていれば、SOCKS宛先を所属回路と経路へ結び付けられる。
コアダンプ取得時点をまたぐ入口PCAPがあれば、TLSレコードからTorセルを取り出し、CGOを処理してリレー本文を回収できる。復元単位は493バイトのTorリレーメッセージ本体である。DATAメッセージにアプリケーションバイトが含まれる場合でも、HTTPSやOnion Service上位プロトコルの暗号文は別途復号しなければならない。

この情報は、侵害端末が調査時点にどの宛先へ接続し、どの回路を使い、どの暗号方式の初期化が確認でき、取得直後にどのリレーメッセージを送ったかを検証する用途に向く。
通常ログに宛先や本文がなく、Control Portへ後から接続できない場合でも、取得済み証拠から再検証できる点が実務上の利点である。

ただし、ローカルホスト側の平文PCAPがすでに取得され、目的がアプリケーション内容の確認だけなら、表~\ref{tab:practical-alternatives}のとおり先にその方法を使うべきである。本手法が追加する価値は、平文らしい断片の発見ではなく、終了後の取得物から所有回路、入口順、ホップ方式、TLSレコード、CGO認証済みリレーメッセージを同じ時点へ帰属させる点にある。実環境へ導入する際は、(1)入口リンクPCAPの継続取得、(2)Tor本体・共有ライブラリ・外部デバッグシンボルの保全、(3)短時間のプロセス停止を伴う整合したコアダンプの取得、(4)秘密ファイルのアクセス制御を取得計画へ含める必要がある。

合成送信元\code{127.0.0.2}と合成識別子を用いた評価では、送信元、ストリーム、宛先、回路ID、3ホップ経路、復号内容中の識別子に関する各検査がすべて一致した。
これは端点メモリ侵害下で匿名性メタデータが失われる仕組みを示す。ただし、実在人物の同定、接続元IPの外部帰属、第三者通信の観測は行っていない。


TLSレコードとCGO層の復号は、CGOの鍵を暗号解析で推定した結果ではない。端点メモリに正規のTLS秘密とCGO状態が存在し、それをTorの所有関係から正しい回路へ帰属させた結果である。
従って脅威モデルは端点メモリ読取りを含む。PCAPだけを持つ受動観測者には適用できない。

本結果は、Tor全体または過去通信の包括的復号を意味しない。本研究が示したのは、取得時点の端点状態を用いた、対象と時間を限定した復号である。


暗号状態を両端で比較する場合、対象プロセスを先に同時停止し、その後でメモリを保存すべきである。順次取得では数秒で状態がずれる。
一方、クライアント単体のコアダンプと同時刻PCAPから、そのコアダンプ以後のセルを処理する場合は、両端メモリの同時取得を必要としない。必要なのは、コアダンプ内のTLS/CGO状態と、PCAP内セルの順序が一致することである。

本実験の成立条件を現場取得で再現するには、パケット取得を継続した状態でTorを停止し、コアダンプ、実行ファイル、外部デバッグシンボル、Build ID、時刻、ハッシュを保存する必要がある。解析では、読み取りの成否と列挙の完了状態を確認する。列挙が不完全な場合、回路の不存在は判定できない。

\section{限界}\label{sec:limitations}

第一に、TLSレコードとCGO層を連続して復号する評価はDebian arm64、自己ビルドTor 0.4.9.16とTor Project公式パッケージ0.4.9.17、OpenSSL 3.0.22、\code{TLS_AES_256_GCM_SHA384}に限られる。x86\_64の構造復元はQEMUゲストコアで確認したが、ネイティブx86\_64でTLSからCGOまでを通す評価は未実施である。OpenSSLのマイナーバージョンやプロバイダーが変われば、内部プロファイルを再検証する必要がある。

第二に、10/10成功は固定条件での反復性を示し、セル・TLS境界を観測しない固定遅延で取得した自己ビルド版6/6試行は、TLSレコード番号と宛先ホップへの一般化を補う。ただし、一般環境での成功率を精密に推定するものではない。交換可能なBernoulli試行と仮定した記述的な95\% Wilsonスコア区間は、それぞれ72.25--100\%、60.97--100\%と広い。公式パッケージ版では一意なCGO認証が6/6で成功し、その区間は60.97--100\%である。一方、評価用マーカーを含むリレーDATAセル本体は5/6であり、その区間は43.65--96.99\%である。追加試行はCGO認証成功ホップ1/2を含むが、ホップ0と異なる暗号スイートは未評価である。また、混在10回路は方式を分類したが、層別通信復元は未評価である。

通信を一時停止した10試行では、Torを\code{SIGSTOP}してからコアダンプを取得した。追加6試行はストリームを止めず、停止時点をセルやTLSレコードの境界に合わせなかったが、\code{gcore}によるメモリ複写中のプロセス停止は残る。実行を止めないライブメモリ取得や、命令途中で不整合なページを収集する取得器は評価していない。したがって本結果は「未改変C Torの制御下の継続通信中でも、整合したスナップショットを取得すれば復元できる」ことを示すが、非原子的なメモリ取得一般を保証しない。

第三に、KeyUpdate、欠損PCAP、TCP部分重複、IP断片化、IPv6、pcapng、後方向、複数連続セルの長時間再生を評価していない。実装のTCP再構成は同一番号の完全再送だけを重複排除する。現行状態から直前セルを復号できなかったが、古い状態の別コピーやメモリアロケータの残留領域が残る可能性を一般には否定できない。

第四に、配布パッケージ版の評価はTor Project公式Debian 12 arm64パッケージ一種類と、同じリポジトリが配布する外部デバッグシンボルに限られる。3/3起動で所有リスト、経路構造、暗号コンテキストを検証した。別の6/6起動ではTLS AEAD認証と一意なCGO認証に成功したが、評価用マーカーを含むリレーDATAセル本体は5/6だった。他のディストリビューション、外部デバッグシンボルを取得できない現場資料、異なるTLS暗号スイートは評価していない。したがって、任意の配布Torへそのまま適用できるとは主張しない。

第五に、C Tor向けの型・所有関係を使用しており、ArtiのRust型へそのまま適用できない。C TorとArtiのCGO逐次状態が正常操作列で一致することは確認したが、ArtiではRustの所有関係に対応する解析器が別途必要である。

第六に、文献調査は主要仕様、査読論文、実装、学位論文を対象としたスコーピングレビューである。未公開ツールや非英語文献を含む系統的レビューは完了していない。このため「世界初」とは主張しない。

最後に、本研究は管理下の隔離網と合成識別子だけを使用した。人物同定、広域の受動監視、遠隔からの匿名性破壊を実証していない。これらを本結果から外挿してはならない。

\section{結論}

本研究は、Torの暗号状態を「鍵らしいバイト列」として探すのではなく、Tor自身の所有関係から回路、入口、ホップ、TLS、CGOをたどるメモリフォレンジック手法を示した。
複数のC Tor 0.4.9系自己ビルドで144/144回路と432/432ホップを復元した。Tor Project公式Debianパッケージ版でも36/36回路と108/108ホップを復元した。

取得時点が対応するクライアントコアダンプとPCAPによる通信復元は、ストリームを一時停止した自己ビルド版10/10独立起動と、セル・TLS境界を観測しない固定遅延で取得した自己ビルド版6/6独立起動で再現した。公式パッケージ版6起動でも、一意なCGO認証済みリレーセル本体を6/6で復元した。このうち評価用マーカーを含むリレーDATAセル本体は5/6だった。残る1起動ではCGO認証に成功した別のリレーセル本体を得たため、目的の通信内容を常に得られるとは主張しない。

通信復元の前提条件は、所有関係、経路構造、暗号コンテキストを検証でき、全ホップでCGOが確立していることである。評価用正解は復号処理から分離し、Control Protocol由来の正解回路IDを候補選択に使わなかった。CGO/\torone{}混在10回路は全ホップCGO条件を満たさず、その通信復元は未評価である。
これにより、保存PCAP、TLSレコード、Tor回路、CGOホップ、493バイトのリレー本文の対応関係を検証できる。

同時に、成立範囲も明確になった。現行CGO状態は過去通信を包括的に開く鍵ではない。直前セルは3/3で復号できず、活動中の順次取得では状態一致が0/36になった。既知の後方向再同期攻撃も、3ホップ処理と回路終了によって、鍵なし復号や匿名性破壊には至らなかった。

一般化には、ネイティブx86\_64、異なるOpenSSLプロバイダー、TLS KeyUpdate、欠損PCAP、後方向、Artiの所有構造に対する評価が残る。拡張時にも、復号範囲だけでなく、説明可能な証拠条件・時点・回路を明示する必要がある。

\printbibliography[title={参考文献}]

\end{document}
